% Document class for LME theses: lmedoc
% %LANGUAGE
% %CONFIG
%    The option "german" uses german.sty
%    For english papers, use the "english" option
% Possible types of theses:
% bt - Bachelor's thesis
% mt - Master's thesis
% diss - Dissertation
% sa - Student thesis
% fp - Forschungspraktikum
% pa - Projektarbeit
%\documentclass[german,bt]{lmedoc/lmedoc}
\documentclass[english,mt]{lmedoc/lmedoc}


%%%%%%%%%%%%%%%%%%
% pdflatex and lualatex are supported
% ++ "Umlaut" support
%    The package "inputenc" can be used to write Umlaute or the german double s
%    directly. You need to use the correct encoding, e.g. latin1.
\usepackage{booktabs}
\usepackage[table]{xcolor}
\usepackage{graphicx} % REQUIRED for resizebox
\usepackage{titlesec}
\usepackage[table]{xcolor}
\usepackage{siunitx}
\usepackage{placeins}
\usepackage{longtable}
\titlespacing*{\section}{0pt}{0pt}{0.5em}
\usepackage{iftex}
\ifPDFTeX
  \usepackage[utf8]{inputenc}
  \usepackage[T1]{fontenc}
  \usepackage{float}
  \usepackage[font=small,labelfont=bf]{caption}
  \usepackage{lmodern}
  \setlength{\parskip}{3pt}
\else
  \ifXeTeX
     \usepackage{fontspec}
  \else 
     \usepackage{luatextra}
  \fi
\fi

%Load additional packages
\usepackage{lmedoc/packages}


% put in this file your abbreviations
% add your abbreviations here

\newabbreviation{ct}{CT}{Computed Tomography}
\newabbreviation{cnns}{CNNs}{Convolutional Neural Networks}
\newabbreviation{GPU}{GPU}{Graphics Processing Unit }
\newabbreviation{mae}{MAE}{Mean Absolute Error}
\newabbreviation{rmse}{RMSE}{Root Mean Squared Error}
\newabbreviation{psnr}{PSNR}{Peak signal-to-noise ratio}
\newabbreviation{mse}{MSE}{Mean Squared Error}
\newabbreviation{ssim}{SSIM}{Structural Similarity Index Measure }
\newabbreviation{flops}{FLOPs}{Floating Point Operations }
\newabbreviation{dicom}{DICOM}{Digital Imaging and Communications in Medicine }
\newabbreviation{hu}{HU}{ Hounsfield Units}
\newabbreviation{DiffCT}{DiffCT}{Differentiable CT framework }
\newabbreviation{sid}{SID}{Source-to-Isocenter Distance }
\newabbreviation{sdd}{SDD}{Source-to-Detector Distance }
\newabbreviation{ROIs}{ROIs}{Regions of Interest}   
\newabbreviation{rlkb}{RLKB}{Residual Large-Kernel  Block }
\newabbreviation{gelu}{GELU}{Gaussian Error Linear Unit }
\newabbreviation{ca}{CA}{Channel Attention}
\newabbreviation{feb}{FEB}{Frequency Enhanced Branch}
\newabbreviation{hpc}{HPC}{High-Performance Computing}
\newabbreviation{mrlkv}{MR-LKV}{Motion Restoration using Large Kernel Convolution }
\newabbreviation{FBP}{FBP}{ Filtered Back Projection}
\newabbreviation{2D }{2D }{Two-dimensional}
\newabbreviation{3D }{3D }{Three-dimensional}
\newabbreviation{MRI}{MRI}{Magnetic Resonance Imaging}
\newabbreviation{1-D}{1-D}{One-dimensional}
\newabbreviation{rf}{RF}{Receptive Field}
\newabbreviation{erf}{ERF}{Effective Receptive Field}
\newabbreviation{xray}{X-Ray}{X-radiation}
\newabbreviation{vif}{VIF}{Visual Information Fidelity}
\newabbreviation{lpips}{LPIPS}{Learned Perceptual Image Patch Similarity }
\newabbreviation{DCGAN}{DCGAN}{Deep Convolutional Generative Adversarial Network}
\newabbreviation{mc-sart}{MC-SART}{Motion-Compensated Simultaneous Algebraic Reconstruction Technique}
\newabbreviation{swinir}{SwinIR}{Swin Transformer for Image Restoration}
\newabbreviation{convnext}{ConvNeXt}{Convolutional Next}
\newabbreviation{slak}{SLaK}{Sparse Large Kernel Network }
\newabbreviation{fft}{FFT}{Fast Fourier Transform}
\newabbreviation{ifft}{IFFT}{Inverse Fast Fourier Transform}

\newabbreviation{fc}{FC}{Fully Connected}
\newabbreviation{ConvFFN }{ConvFFN }{Convolutional Feed-Forward Network }
\newabbreviation{rstbs}{RSTBs}{Residual Swin Transformer Blocks }
\newabbreviation{mdta}{MDTA}{Multi-DConv Head Transposed Attention }
\newabbreviation{cbct}{CBCT}{Cone-beam Computed Tomography  }
\newabbreviation{fdk}{FDK}{Feldkamp–Davis–Kress}


%%%%%%%%%%%%%%%%%%%%%%%%%%%%

% supress messages because of underful hbox, this is not a problem
\hbadness=10000



% some useful commands
\makeatletter % let's define a single dot
\DeclareRobustCommand\onedot{\futurelet\@let@token\@onedot}
\newcommand{\@onedot}{\ifx\@let@token.\else.\null\fi\xspace}
\makeatother

% Sets the bib file

\addbibresource{literature.bib}

% When writing a large document, it is sometimes useful to work on selected sections of the document.
% Use this command to only build the document partially. Speeds up the development cycle.
% For the final product, this has to be commented out.
%\includeonly{introduction, appendix, foo, bar}


\pagenumbering{roman}

\begin{document}
\clearpage
% %CONFIG
% This is for students' theses
  \begin{deckblatt}
    \Titel{Large Kernel Convolution for CT Image Restoration} % Title
    \Name{Bhattacharjya} % Last name
    \Vorname{Debanjona} % Given name
    \Geburtsort{Siliguri, India} % Place of birth
    \Geburtsdatum{16.11.1998} % Date of birth
    \Betreuer{Yipeng Sun, Prof Dr. Ing Andreas Maier} % Advisor
    \Start{17.10.2025} % Start of thesis
    \Ende{16.04.2026} % End of thesis
    %\ZweitInstitut{ZweitInstitut} % Cooperation partner
  \end{deckblatt}


\cleardoublepage
% Only necessary for thesis

Ich versichere, dass ich die Arbeit ohne fremde Hilfe und ohne Benutzung
anderer als der angegebenen Quellen angefertigt habe und dass die Arbeit
in gleicher oder "ahnlicher Form noch keiner anderen Pr"ufungsbeh"orde
vorgelegen hat und von dieser als Teil einer Pr"ufungsleistung
angenommen wurde. Alle Ausf"uhrungen, die w"ortlich oder sinngem"a"s
"ubernommen wurden, sind als solche gekennzeichnet.
\\


\cleardoublepage


\begin{center}
\bfseries
% Abstract in German
{\selectlanguage{english}"Ubersicht}
\normalfont
\end{center}


\vspace{5.0cm}

\begin{center}
\bfseries
% Abstract in English
{\selectlanguage{english}Abstract}
\normalfont
\end{center}
\vspace{1cm}
Motion artifacts in computed tomography (CT) significantly compromise image quality and can adversely affect diagnostic accuracy. These artifacts arise from inconsistencies in projection data caused by patient motion and often exhibit complex global structures that are challenging to correct using conventional convolutional neural networks (CNNs) with limited receptive fields. This thesis introduces a deep learning framework based on a super‑large-kernel convolutional network (MR‑LKV) designed to effectively reduce motion artifacts. The model employs super-large-kernel convolutions to efficiently capture long-range spatial dependencies, enabling effective modeling of globally distributed inconsistencies. The framework is examined in both the projection domain and the image domain: in the projection domain, the model restores corrupted sinograms by correcting view‑wise inconsistencies, while in the image domain, it refines reconstructed CT volumes affected by motion. Experimental evaluations demonstrate that MR‑LKV achieves strong performance across both domains. In the projection domain, the model effectively mitigates the source of motion artifacts, resulting in superior reconstruction quality. In the image domain, it attains competitive visual results relative to CNN‑ and transformer‑based baselines while maintaining substantially lower computational cost and faster inference. Overall, the findings indicate the effectiveness of large‑kernel CNNs as a powerful and efficient approach for modeling global spatial dependencies in CT reconstruction, offering a practical and scalable solution for motion artifact correction.

 

 
\clearpage
\begin{center}
\bfseries
% Abstract im Deutsch
{\selectlanguage{english}Zusammenfassung}
\normalfont
\end{center}
\vspace{1cm}
Bewegungsartefakte in der Computertomographie (CT) beeinträchtigen die Bildqualität erheblich und können die diagnostische Genauigkeit negativ beeinflussen. Diese Artefakte entstehen durch Inkonsistenzen in den Projektdaten infolge von Patientenbewegungen und weisen häufig komplexe globale Strukturen auf, die mit konventionellen Convolutional Neural Networks (CNNs) mit begrenztem Rezeptionsfeld nur schwer korrigiert werden können. Diese Arbeit stellt ein Deep-Learning-Framework basierend auf einem Convolutional Neural Network mit sehr großen Faltungsfiltern (MR-LKV) vor, das speziell zur Reduktion von Bewegungsartefakten entwickelt wurde. Das Modell nutzt supergroße Faltungskerne, um langfristige räumliche Abhängigkeiten effizient zu erfassen und somit global verteilte Inkonsistenzen wirksam zu modellieren. Der Ansatz wird sowohl im Projektions- als auch im Bildraum untersucht: Im Projektionsraum rekonstruiert das Modell gestörte Sinogramme durch die Korrektur von Inkonsistenzen zwischen einzelnen Projektionen, während es im Bildraum rekonstruierte CT-Volumina mit Bewegungsartefakten direkt verbessert. Experimentelle Ergebnisse zeigen, dass MR-LKV in beiden Domänen eine starke Leistung erzielt. Im Projektionsraum werden die zugrunde liegenden Ursachen der Bewegungsartefakte effektiv reduziert, was zu einer deutlich verbesserten Rekonstruktionsqualität führt. Im Bildraum erreicht das Modell wettbewerbsfähige visuelle Ergebnisse im Vergleich zu CNN- und transformerbasierten Ansätzen, bei gleichzeitig deutlich geringerer Rechenkomplexität und kürzeren Inferenzzeiten. Insgesamt zeigen die Ergebnisse, dass großkernige CNN-Architekturen einen effektiven und effizienten Ansatz zur Modellierung globaler räumlicher Abhängigkeiten in der CT-Rekonstruktion darstellen und somit eine praktische und skalierbare Lösung zur Reduktion von Bewegungsartefakten bieten.

 
\cleardoublepage

\tableofcontents

\cleardoublepage \pagenumbering{arabic}

\chapter{Introduction}


\noindent

Computed tomography (CT) has become a vital tool of modern medical diagnosis due to its ability to provide high-resolution cross-sectional images of internal anatomical structures. It plays a crucial role in emergency diagnostics, neurological assessment, trauma assessment, and treatment planning. In many diagnostic applications, clinicians rely on fine structural details, placing strong requirements on the quality and reliability of reconstructed CT images. Even slight degradations in image quality can negatively affect diagnostic confidence and increase the risk of misinterpretation ~\cite{common_ct_artifacts, kak2001principles}. 


Even with major advancements in CT hardware and reconstruction methods, image quality remains susceptible to degradation from a range of acquisition-related artifacts ~\cite{common_ct_artifacts}. Among these, patient motion continues to be one of the most frequent and challenging causes of degradation. Motion may arise from involuntary physiological processes such as breathing and cardiac activity, or from a patient’s inability to remain completely still during the scan ~\cite{boas_ct_artifacts}. Standard CT reconstruction techniques work under the assumption that the object remains still throughout data acquisition, meaning all projection angles capture the same anatomy. When motion occurs, this assumption breaks down and as a consequence, the projections no longer align with a single consistent structure  ~\cite{common_ct_artifacts}. Consequently, the reconstruction task becomes ill-posed, as no single static image can accurately represent the motion-affected data~. These inconsistencies are carried through the reconstruction pipeline, manifesting as visible artifacts such as streaking, blurring, or even repeated structures in the final images ~\cite{boas_ct_artifacts}. Such effects can obscure important anatomical details and compromise diagnostic reliability.
An example of motion-induced degradation is shown in Fig. \ref{fig:cleanVsArtifacted}, where streaking artifacts and loss of structural detail are clearly visible compared to the clean reconstruction. 

\begin{figure}[htpb]
    \centering
    \includegraphics[width=0.5\linewidth]{cleanvsartifactedroi.png}
    \caption{Visual comparison of clean and motion-corrupted CT reconstructions. Zoomed regions (ROIs) emphasize the loss of fine anatomical details and reduced edge sharpness}
    \label{fig:cleanVsArtifacted}
\end{figure}

Conventional motion compensation approaches in CT are commonly formulated as a two-step pipeline consisting of motion estimation followed by motion-compensated reconstruction, often based on optimization strategies and imaging models  ~\cite{thies2025_cbctmotion}. These approaches are effective when motion can be described using simplified or well-constrained models. However, in practice, patient motion can be complex and difficult to estimate accurately, particularly in the presence of irregular or non-rigid motion patterns ~\cite{boas_ct_artifacts}. This limits the robustness of purely model-based correction techniques and can lead to performance degradation under realistic conditions. As a result, such methods tend to be fragile outside controlled settings, which has encouraged the search for data-driven alternatives that can learn complex artifact patterns without explicit motion modeling ~\cite{datadriven23}. 

In recent years, deep learning models have shown remarkably strong results for medical image reconstruction and restoration ~\cite{ahishakiye2021survey, modality_dl_recon}. Rather than depending on hand-crafted correction models, these methods learn intricate relationships between degraded measurements and their corresponding high-quality reconstructions directly from training data. Convolutional neural networks (CNNs) have been particularly successful in image restoration tasks due to their ability to learn hierarchical spatial features and represent intricate degradation patterns that are challenging to describe analytically ~\cite{Goodfellow2016, modality_dl_recon}. These successes have sparked growing interest in applying deep learning techniques to motion artifact reduction in CT imaging. 

However, most existing CNN-based approaches rely on small convolution kernels, such as 3×3 or 5×5, and expand the receptive field primarily by stacking many layers ~\cite{simonyan2015vgg}.  While this design is effective for a wide range of image restoration problems, it may be suboptimal for representing motion artifacts in CT. These artifacts can span large spatial extents and give rise to long-range streaks in the reconstructed slices ~\cite{boas_ct_artifacts}. Modeling such globally distributed structures with only small kernels typically requires very deep networks and intricate multi-scale schemes, which in turn increases computational burden and reduces efficiency.


Recent progress in CNN architecture design has renewed the use of large convolutional kernels to capture long-range spatial dependencies ~\cite{replknet2022, Luo2016}. Enlarging the kernel size allows individual layers to access a wider spatial context, which makes it easier to model degradations that are distributed over large regions. At the same time, modern hardware and architectural optimization techniques have made it feasible to implement large-kernel CNNs efficiently ~\cite{replknet2022}. Altogether, these developments suggest that super-large-kernel convolution offers a promising approach for addressing motion artifacts in CT imaging.


Transformer-based architectures have recently become popular in image restoration due to their ability to model global context through self-attention mechanisms ~\cite{vaswani2017attention,restormer2022}. While they often achieve strong performance, they typically require substantial computational resources and memory. In contrast, CNNs benefit from highly optimized implementations and broad hardware support. For high‑resolution CT data and potential real‑time clinical use, CNNs with improved long‑range modeling offer a practical and computationally efficient alternative ~\cite{replknet2022}.


\section{Research Gap}

The preceding discussion highlights that motion artifacts in CT images are inherently global phenomena, often extending across large spatial regions. Although many CNN-based restoration approaches have shown good performance, they generally employ small convolution kernels and enlarge the receptive field mainly by increasing network depth, which can be inefficient for representing such spatially extended artifacts ~\cite{simonyan2015vgg, Luo2016}. 

While transformer-based models effectively capture global contextual information, their high computational and memory requirements make them less suitable for large-scale CT imaging scenarios ~\cite{restormer2022, liang2021swinir}. This highlights the need for alternative approaches that can capture long-range dependencies while maintaining computational efficiency. 
Specifically, this thesis addresses this gap by investigating whether super-large-kernel convolutional networks, applied on both the projection domain and image domain, can capture long-range spatial dependencies in CT sinogram data and correct motion-induced inconsistencies before they propagate through the reconstruction process, while maintaining the computational efficiency of convolutional architectures. 
Fig.\ref{fig:research_gap}  shows the research gap that is being addressed in this thesis.

\begin{figure}
    \centering
    \includegraphics[width=1.0\linewidth]{research gap.jpg}
    \caption{Illustration of the research gap in CT artifact correction}
 
    \label{fig:research_gap}
\end{figure}

\section{Objective}


The objective of this thesis is to investigate the effectiveness of super-large-kernel convolutions for reducing motion artifacts in CT imaging by directly modeling long-range spatial dependencies. 
To this end, this work develops and evaluates a CNN architecture incorporating large convolutional kernels for artifact reduction, assessing its performance against conventional small-kernel approaches.\

The proposed model, referred to as MR-LKV, is designed to capture extended artifact structures while maintaining computational efficiency. Its performance is benchmarked against several baseline architectures, including traditional convolutional networks and transformer-based restoration models.
\clearpage
\section{Contributions}
The main contributions of this thesis are summarized as follows:

• Development of projection-domain and image-domain restoration frameworks for CT motion artifact reduction,  combining forward projection, controlled motion simulation, and deep-learning-based CT image restoration.

• Introduction of MR-LKV, a novel convolutional architecture with super-large-kernel convolutions designed to capture long-range spatial dependencies in CT data.


• Incorporation of frequency-enhanced feature integration and view-axis attention mechanisms to improve artifact modeling within projection space.


• Comprehensive experimental comparison of the proposed approaches against CNN- and transformer-based restoration models in terms of reconstruction quality and computational efficiency.

\section{Thesis Organization}

The remainder of this thesis is organized as follows.

\textbf{Chapter 2 }reviews the theoretical background, including CT fundamentals, the origin of motion artifacts, and deep learning concepts relevant to image restoration. 

\textbf{Chapter 3} reviews existing approaches for CT image restoration, including image-domain methods, model-based reconstruction techniques, and projection-domain learning strategies. 

\textbf{Chapter 4} explores the novel architecture as well as baseline architectures in detail. 

\textbf{Chapter 5 }describes the proposed methodology, covering data preparation, training strategy, inference strategy, and evaluation method. 

\textbf{Chapter 6} presents the experimental results and comparative evaluation. 

\textbf{Chapter 7} summarizes the thesis and outlines potential directions for future work. 





   % Introduction 
\cleardoublepage
\chapter{Theoretical Background}
\label{chap:theoretical_background}
\section{Computed Tomography (CT) Fundamentals}

Computed tomography (CT) is a medical imaging technique in which X-ray projections are received from multiple angles around the patient and later combined to reconstruct cross-sectional images of the body   ~\cite{kak2001principles}. In modern CT scanners, an X-ray source and detector array rotate around the patient, acquiring projection measurements from multiple angular positions. In practice, this process relies on the interaction of X-ray photons with human tissue during image acquisition. As X-ray beams traverse the body, they are attenuated to different degrees depending on the density and composition of the encountered tissues  ~\cite{kak2001principles}.
The attenuation of X-ray intensity along a ray path can be described by the Beer–Lambert law:

\begin{equation}
I = I_{0} \exp\left(- \int_{L} \mu(x)\,dx \right)
\end{equation}

where $I_0$ is the initial intensity of the ray path $L$, $\mu(x)$  is the linear attenuation coefficient at position $x$, and the integral $\int_{L} \mu(x)\,dx$ represents the cumulative attenuation along that path ~\cite{Buzug}. By acquiring such measurements from many projection angles and applying suitable reconstruction algorithms, the spatial distribution $\mu$ can be estimated to produce tomographic images. These projection measurements represent line integrals of the attenuation distribution along different ray paths through the object. In clinical CT systems, the set of line integrals (projection data) measured along all projection angles and detector positions is often represented as a \textit{sinogram}. 
Fig.\ref{fig:CT scan machine} shows the standard overview of a CT scanning process.
\begin{figure}
    \centering
    \includegraphics[width=1.0\linewidth]{ctscan.jpg}
    \caption{Overview of a standard CT scanning process from image acquisition to Filtered Back Projection (FBP). Concept adapted from standard CT imaging literature ~\cite{kak2001principles}}
    \label{fig:CT scan machine}
\end{figure}

\section{Sinograms and Projection Data Representation}

In computed tomography, the raw measurements acquired by the scanner correspond to projection data obtained from multiple viewing angles around the object. Each projection records how much the X‑rays are attenuated along a bundle of rays that pass through the object and hit a row of detector elements. Each detector element, therefore, measures the total attenuation of X-rays along a specific ray path through the object. If these projections are stacked according to their acquisition angle, they form what is known as a \textit{sinogram} \cite{Buzug}. 

Each row of a sinogram corresponds to a projection acquired at a specific angle, while each column tracks a specific detector element over all angles. As the gantry rotates, a single point in the object traces out a sinusoid-like curve in this 2D representation, which is the reason for the name \textit{sinogram}  \cite{Buzug, kak2001principles}. 

Mathematically, the projection measurements acquired by a CT scanner can be described using the Radon transform~\cite{natterer_ct}. For an object with attenuation distribution $f(x,y)$, the projection measurement at angle $\theta$ and detector position $t$ can be written in the following equivalent form: 

\begin{equation}
p(t,\theta) = \int f(x,y)\,\delta(x\cos\theta + y\sin\theta - t)\,dx\,dy ,
\end{equation}

where $p(t,\theta)$ denotes the sinogram value at projection angle $\theta$ and detector position $t$ ~\cite{natterer_ct, kak2001principles}. Reconstruction algorithms such as filtered backprojection aim to recover the attenuation distribution $f(x,y)$ from these projection measurements.
This mapping is illustrated in Fig.\ref{fig:sinogram_formation_in_CT}
 \begin{figure}[H]
     \centering
     \includegraphics[width=0.8
     \linewidth]{Presentation2.jpg}
     \caption{Illustration of sinogram formation in CT, where the object $f(x,y)$ is mapped to projection data $p(t,\theta)$ via the Radon transform and reconstructed by backprojection. Adapted from ~\cite{kak2001principles}}
     \label{fig:sinogram_formation_in_CT}
 \end{figure}
Sinograms provide an intermediate representation that is directly related to the acquisition process of CT scanners. Reconstruction algorithms implicitly assume that all projections correspond to a \textbf{single static object configuration} ~\cite{common_ct_artifacts}. If this assumption is violated, for example, due to patient motion during the scan, the sinogram no longer represents a single, consistent object state ~\cite{common_ct_artifacts}. These inconsistencies manifest as distortions or discontinuities in the sinogram and propagate into structured artifacts in the reconstructed image. An example sinogram is shown in Fig. \ref{fig:clean_vs_artifacted_sinogram}. 

 \begin{figure}[H]
     \centering
     \includegraphics[width=0.95\linewidth]{sinogram.jpg}
     \caption{Example sinogram obtained from CT projection measurements. Each row corresponds to a projection acquired at a different projection angle, while each column represents a detector position. The sinusoidal patterns correspond to the trajectories of object points across different projection angles. (left): Clean Sinogram; (right): Motion-corrupted Sinogram}
     \label{fig:clean_vs_artifacted_sinogram}
 \end{figure}

\section{CT Acquisition Geometry}
Depending on the configuration of the X-ray source and detector, different geometries are used in practice, including parallel-beam, fan-beam, and cone-beam geometries. 

\subsection{Parallel-Beam Geometry}

In parallel-beam geometry, all X-ray rays are assumed to be parallel to each other as they pass through the object. Each projection represents line integrals taken along parallel paths at different angles \cite{kak2001principles}.  Although parallel-beam geometry is not directly used in modern clinical CT systems, it provides the theoretical foundation for many reconstruction algorithms, including the classical Radon transform and filtered backprojection.

\subsection{Fan-Beam Geometry}

Fan-beam geometry, widely used in conventional clinical CT scanners, uses X-rays that diverge from a point source in a two-dimensional plane and are detected by a one-dimensional detector array \cite{Buzug}. Compared to parallel-beam geometry, fan-beam acquisition is more efficient and better suited for practical systems, although it introduces additional geometric complexity that must be accounted for during reconstruction.

\subsection{Cone-Beam Geometry}

Cone-beam computed tomography (CBCT) uses a divergent cone-shaped X-ray beam and a flat-panel detector to capture volumetric data in one rotational scan \cite{feldkamp1984practical}. Unlike fan-beam CT, which reconstructs sequential slices, CBCT generates a full 3D volume by rotating around the object.

Each CBCT projection is a 2D image, and together they form a 3D dataset spanning multiple viewing angles. This geometry enables efficient volume reconstruction but also increases sensitivity to noise, scatter, and motion.

Reconstruction is commonly done using the Feldkamp–Davis–Kress (FDK) algorithm \cite{feldkamp1984practical}, an extension of filtered backprojection for cone-beam geometry. However, due to the divergent beam and circular trajectory, FDK provides only an approximate reconstruction and is highly affected by projection inconsistencies.

The fundamental differences between CT acquisition geometries, including parallel-beam, fan-beam, and cone-beam configurations, are illustrated in Figure \ref{fig:ct_geometry}.

\begin{figure}[htpb]
    \centering
    \includegraphics[width=1.0\linewidth, height=0.2\textheight ]{ct_geometry.pdf}
    \caption{Comparison of CT acquisition geometries.  }
    \label{fig:ct_geometry}
\end{figure}

\section{CT Reconstruction Algorithms}
Once projection measurements are collected, the internal attenuation distribution of the scanned object must be reconstructed from the sinogram data. This reconstruction process aims to recover the spatial distribution of linear attenuation coefficients from projection data. Mathematically, it involves inverting the \textbf{Radon} transform that relates these projections to the object's underlying structure. 
\subsection{Filtered Back Projection}
Filtered Back Projection (FBP) is a common reconstruction method in clinical CT systems ~\cite{Buzug}. In simple backprojection, each projection value is distributed back along the corresponding ray paths across the image domain. This alone gives blurred images due to overemphasis on low-frequency components. Hence, FBP applies a high-pass filter to the projections in the frequency domain before backprojecting across all angles, enabling efficient recovery of the attenuation map. Due to its computational efficiency and well-understood behavior, FBP remains widely used in many CT imaging systems.

However, FBP is highly sensitive to inconsistencies in projection data. Even small deviations from the assumption of a static object, such as motion, can lead to significant artifacts in the reconstructed image.

\subsection{Iterative Reconstruction}
Beyond analytical methods like FBP, modern CT systems often use iterative reconstruction techniques. These methods formulate CT reconstruction as an optimization problem and iteratively refine an estimate of the image by minimizing the difference between measured and simulated projections ~\cite{Buzug}. While offering better noise reduction and artifact handling, iterative methods require substantially more computation than analytical approaches. In the presence of complex or unpredictable motion, their performance may degrade.

\subsection{Feldkamp–Davis–Kress}
The Feldkamp–Davis–Kress (FDK) \cite{feldkamp1984practical} algorithm extends the filtered backprojection method for cone-beam CT systems, where projections are acquired from a divergent X‑ray source onto a 2D detector. To reconstruct the 3D attenuation map, FDK applies three key steps: geometric weighting to correct beam divergence, frequency-domain filtering of each projection with a ramp filter (optionally windowed to suppress noise), and cone-beam backprojection over all angles. Although only an approximate solution, it remains the standard for circular scanning due to its efficiency. Like conventional backprojection, FDK is prone to artifacts from inconsistent data, including those caused by patient motion. 



 

\section{Motion Artifacts in Computed Tomography (CT)}

Although reconstruction algorithms aim to recover the attenuation distribution from projection measurements, they generally assume that the scanned object remains stationary during data acquisition ~\cite{common_ct_artifacts}.  In clinical practice, patient motion can occur due to involuntary movements such as breathing, physiological activity, or the inability to remain completely still during the scan ~\cite{boas_ct_artifacts}. When motion occurs during acquisition, projections measured at different angles no longer correspond to the same object configuration. As a result, inconsistencies arise in the projection data. Since reconstruction algorithms combine all projections under the assumption of a static object, these inconsistencies propagate through the reconstructed images and manifest as visible artifacts such as streaks or blurs  ~\cite{common_ct_artifacts}.  Motion artifacts can conceal important anatomical details and significantly reduce diagnostic reliability. 
Fig \ref{fig:cleanVsartifactedctimage} shows a brief overview of affect of motion artifact on head CT scan. 

\paragraph{Types of Motion in CT Acquisition}

Motion during CT acquisition can be broadly categorized based on its nature and temporal behavior:

\begin{itemize}

\item \textbf{Rigid motion:}  
This involves the entire object translating or rotating as a whole, such as head movements \cite{ko2021rigid}. It causes projection misalignment and typically produces structured streak artifacts in the reconstruction.


\item \textbf{Non-rigid motion:}  
This includes deformation of anatomical structures, such as breathing or cardiac motion \cite{ko2021rigid}. This leads to spatially varying inconsistencies that are harder to model and often produce complex blurring and distortion.

\item \textbf{Periodic motion:}  
Occurs due to physiological processes such as respiration or heartbeat \cite{common_ct_artifacts}. These motions introduce cyclic inconsistencies in projection data, often resulting in ghosting or duplication artifacts.

\item \textbf{Sudden or abrupt motion:}  
Short-duration movements during acquisition cause localized inconsistencies, leading to severe streak artifacts and discontinuities in the reconstructed image \cite{common_ct_artifacts}.

\end{itemize}




\begin{figure}[htpb]
    \centering

    \begin{minipage}{0.48\textwidth}
        \centering
        \includegraphics[width=\textwidth]{CQ500CT19 CQ500CT19_mid.png}
    \end{minipage}
    \hfill
    \begin{minipage}{0.48\textwidth}
        \centering
        \includegraphics[width=\textwidth]{CQ500CT19 CQ500CT19_artifact_mid.png}
        
    \end{minipage}

    \caption{Illustration of a motion-free axial head CT slice  (left) and the corresponding motion-artifacted axial head CT slice  (right).}
    \label{fig:cleanVsartifactedctimage}
\end{figure}


\section{Deep Learning for Medical Image Restoration}

Medical image restoration aims to improve image quality by reducing artifacts introduced during acquisition and reconstruction. Classical restoration approaches often rely on analytical models or handcrafted regularization techniques. While these methods can be effective under simplified assumptions, they often struggle to model complex and non-linear degradation patterns observed in real clinical imaging scenarios.

In contrast, deep learning methods learn mappings from corrupted to high-quality images, directly from data, allowing them to capture intricate relationships between measurements, reconstruction artifacts, and underlying anatomy ~\cite{ahishakiye2021survey, yedder2021survey}.  In medical imaging, supervised deep learning frameworks commonly train neural networks on paired degraded inputs and high‑quality reference images, for example, in accelerated MRI and CT reconstruction and denoising ~\cite{modality_dl_recon}.  These approaches leverage large datasets to learn complex spatial patterns associated with different types of image degradation.
In the context of CT imaging, both image-domain and projection-domain learning approaches have been explored. 
The general framework is depicted in Fig. \ref{fig:artifacted_to_clean_output}. 

Among deep learning models, convolutional neural networks (CNNs) have become the dominant architecture for image restoration tasks due to their ability to efficiently capture spatial correlations in image data. The following section, therefore, introduces the fundamental principles of CNN architectures used in medical image restoration.

\begin{figure}[H]
     \centering
     \includegraphics[width=0.8\linewidth]{deep learning.jpg}
     \caption{Deep learning-based image restoration framework. A neural network learns a mapping $f_\theta(x)$ from corrupted input images to high-quality restored outputs. }
     \label{fig:artifacted_to_clean_output}
 \end{figure}

\section{Convolutional Neural Networks}

Convolutional Neural Networks (CNNs) are deep learning architectures designed for image analysis and object recognition tasks~\cite{sarvamangala2021cnnmed,cnn_medimage_review}.  They use convolutional layers with learnable filters that operate across the spatial dimensions of an image to extract meaningful features while preserving local spatial structure.
\begin{equation}
y(i,j) = \sum_{m}\sum_{n} x(i+m, j+n)\, w(m,n)
\end{equation}
where ${x}$ is input image, \textbf{$w$} is convolution kernel and $y$ is output feature map \cite{Goodfellow2016}.

By stacking multiple layers, CNNs build a hierarchical representation of the input data. Early layers typically focus on low-level patterns such as edges or textures. As data moves through the network, later layers combine the low-level patterns detected in early layers to form more abstract representations ~\cite{simonyan2015vgg}. 

A generic CNN restoration architecture is shown in Fig. \ref{fig:CNN_architecture} 
\begin{figure}[htbp]
    \centering
    \includegraphics[width=0.7\linewidth]{cnn image.jpg}
    \caption{Generic convolutional neural network for image restoration }
    \label{fig:CNN_architecture}
\end{figure}
This progressive abstraction makes CNNs particularly well-suited for medical image restoration tasks, where both fine structural details and broader contextual information are important. 






\section{Receptive Fields and Kernel Size in CNNs}
The receptive field of a CNN refers to the spatial region of the input that can influence a particular output activation ~\cite{Luo2016}. The theoretical receptive field (RF) of a unit in a convolutional neural network (CNN) defines the spatial extent of the input that can influence its output value. As layers are stacked, this field grows according to a recursive relationship governed by the kernel size, stride, and dilation of each layer ~\cite{Luo2016}.  As a result, deeper networks are, in principle, able to aggregate information from increasingly larger spatial contexts. For example, for a 1-D sequence of layers with kernel size $k_l$ and stride $s_l$, the receptive field $R_l$ at layer $l$  can be computed recursively as 

                 \begin{equation}
R_l = R_{l-1} + (k_l - 1)\prod_{i=1}^{l-1} s_i
\end{equation}
which illustrates how stacking many layers with small kernels can still lead to large receptive fields. ~\cite{Luo2016}.

Although the theoretical receptive field grows with network depth, the region that actually contributes to a network’s response is often smaller. This region is referred to as the \textit{effective receptive field (ERF)}. Empirical studies have shown that the ERF tends to concentrate near the center of the theoretical receptive field and decreases rapidly toward the boundaries  ~\cite{Luo2016}.  Consequently, networks that depend exclusively on small convolution kernels often need numerous layers to capture large-scale image structures effectively, thereby increasing model depth and computational complexity.

This distinction between theoretical receptive field and effective receptive field is illustrated in Fig. \ref{fig:receptive_field}. 

 
 Conventional CNN architectures in image restoration, therefore, often employ small kernels, such as $3\times3$, combined with depth, pooling, or multiscale structures to expand the receptive field ~\cite{simonyan2015vgg}. Such deep stacks increase network depth and computational complexity while still relying on indirect accumulation of contextual information.

 

 
 Motion artifacts in CT images, however, frequently appear as spatially extended patterns, such as streaks and blurring that span large portions of the reconstruction rather than strictly local perturbations~\cite{common_ct_artifacts}. As a result, networks designed for CT motion artifact reduction must access long‑range contextual information to distinguish true anatomical structures from motion‑induced intensity inconsistencies and to correct artifacts that propagate across wide image regions. 

 This limitation becomes particularly critical in CT motion artifact correction, where distortions frequently span large spatial regions. Modeling such global structures using only small kernels is inefficient and often requires deep architectures with increased computational complexity.

This mismatch between local operations and global artifact structure motivates the need for architectures that can directly capture long-range dependencies ~\cite{replknet2022}.

 
\begin{figure}[htbp]
    \centering
    \includegraphics[width=0.6\linewidth]{Receptive Field.png}
    \caption{Conceptual illustration of receptive field expansion. Stacked small convolution kernels gradually increase the receptive field through multiple layers, whereas a single large convolution kernel can directly capture a wider spatial context within one layer. }
    \label{fig:receptive_field}
\end{figure}

\section{Super-Large-Kernel Convolution}

    Super-large-kernel convolution uses convolutional filters with substantially larger spatial dimensions than those typically used in standard CNN architectures, such as 3 × 3 or 5 × 5. By increasing the kernel size, a single convolutional layer can directly capture a wide spatial context, allowing explicit modeling of long-range dependencies without using very deep stacks of small-kernel layers. 
Recent work on large-kernel CNNs has demonstrated that such approaches can be implemented efficiently without prohibitive computational costs~\cite{replknet2022, PeLK2024}. Super-large kernel convolutions provide a direct mechanism for integrating information over the large regions within a single layer or block, potentially reducing the need for very deep small-kernel networks or heavy self-attention modules~\cite{replknet2022}. 

Fig. \ref{fig:super_large_kernels} compares the spatial coverage of small and large kernels. 

One important development is the use of depthwise separable convolutions, in which spatial filtering is performed independently for each channel and followed by pointwise ($1 \times 1$) convolutions for channel mixing. This significantly reduces computational complexity compared to standard convolutions while still allowing the use of very large spatial kernels. Architectures such as RepLKNet ~\cite{replknet2022} have demonstrated that convolution kernels as large as $31 \times 31$  can be employed efficiently and can improve performance on various vision tasks by capturing broader spatial context. 

Combined with modern GPU hardware and efficient implementations, these architectures can be implemented with manageable computational and memory requirements, motivating their use in CT motion artifact reduction tasks where both global contextual awareness and computational efficiency are essential. 
For CT motion artifact correction, this property offers particular advantages since motion-induced inconsistencies typically span multiple projection views in the sinogram. Similarly, in the image domain, motion artifacts appear as globally distributed streaks and blurring, which also benefit from large receptive field modeling. Large convolution kernels, therefore, provide a natural mechanism for modeling these extended artifact structures. 

\begin{figure}[htbp]
    \centering
    \includegraphics[width=1.0
    \linewidth]{large kernel.jpg}
    \caption{Comparison between small and large convolution kernels. Small kernels capture limited local context, whereas large kernels significantly increase the receptive field, enabling the network to model global spatial dependencies.}
    \label{fig:super_large_kernels}
\end{figure}

\section{ CNN versus Transformers for Image Restoration}

Transformer-based architectures have recently gained prominence in image processing due to their ability to model long-range dependencies through self-attention mechanisms \cite{vaswani2017attention}. In image restoration tasks, transformer models can effectively capture global context, which is beneficial for handling structured image degradations ~\cite{restormer2022}. 
Despite their strengths, transformer models often incur higher computational complexity and memory consumption compared to convolutional architectures. These factors can pose challenges for deployment in resource-constrained or real-time medical imaging environments. 
At the same time, recent studies indicate that appropriately designed CNNs, particularly those with enlarged receptive fields, can close much of the performance gap compared to transformer-based models on various image restoration tasks, while saving on computation time~\cite{replknet2022}. Given the high spatial resolution and real-time constraints often present in CT imaging, efficient CNN architectures with strong global modeling capabilities provide a practical and competitive solution for motion artifact reduction.

A summary of these differences is provided in Table \ref{tab:attention_vs_lk} 
\begin{table}[H]
\centering
\caption{Comparison between self-attention and large-kernel depthwise convolution}
\label{tab:attention_vs_lk}
\renewcommand{\arraystretch}{1.2}
\begin{tabular}{lcc}
\toprule
\textbf{Property} & \textbf{Self-Attention} & \textbf{Large-Kernel Depthwise Conv} \\
\midrule
Complexity per layer & $\mathcal{O}(n^2 d)$ & $\mathcal{O}(n k^2 c)$ \\
Scaling with resolution & Quadratic & Linear \\
Receptive field & Global (full input) & Local (kernel size $k$) \\
Parameter count per layer & $\mathcal{O}(d^2)$ & $\mathcal{O}(k^2 c)$ \\
Hardware optimization & Moderate & Strong (GPU-optimized) \\
Memory footprint & High (attention maps) & Low (no attention maps) \\
\bottomrule
\end{tabular}
\end{table}

\section{Image Quality Metrics and Loss Functions}

The performance of image restoration models is evaluated using a combination of pixel-wise, structural, perceptual, and computational efficiency metrics. This multifaceted evaluation ensures a comprehensive assessment of reconstruction quality and practical applicability.

\subsection{Pixel-wise Metrics}

The\textbf{ L1 loss}, also known as the \textbf{mean absolute error (MAE)}, measures the average absolute difference between predicted and reference values. For a predicted image $\hat{\mathbf{S}}$ and ground-truth image $\mathbf{S}$, the L1 loss is defined as

\begin{equation}
\mathrm{MAE} 
= \frac{1}{N} \sum_{i=1}^{N} 
\left| \hat{S}_i - S_i \right|,
\end{equation}

where $N$ denotes the number of pixels (or sinogram elements). Unlike  L2 loss, L1 loss is less sensitive to outliers and tends to preserve sharper edges, making it widely used in image restoration tasks. 



The \textbf{Root Mean Squared Error (RMSE) } is a widely used metric for quantifying reconstruction accuracy by measuring the average magnitude of squared differences between predicted and reference values. It is defined as:

\begin{equation}
\mathrm{RMSE} = \sqrt{\frac{1}{N} \sum_{i=1}^{N} \left(\hat{S}_i - S_i\right)^2}
\end{equation}

where $\hat{S}_i$ and $S_i$ denote the predicted and ground truth values, respectively, and $N$ represents the total number of pixels.

RMSE penalizes larger errors more strongly than the Mean Absolute Error (MAE) due to the squaring operation, making it particularly sensitive to outliers and significant reconstruction deviations. As a result, RMSE provides complementary information to MAE by emphasizing regions with high reconstruction error.



\textbf{Peak Signal-to-Noise Ratio (PSNR}) is a widely used metric for measuring reconstruction quality. It is derived from the mean squared error (MSE) between a reconstructed image and its reference:

\begin{equation}
\mathrm{PSNR} = 10 \log_{10} 
\left( \frac{L^2}{\mathrm{MSE}} \right),
\end{equation}

where $L$ denotes the maximum possible pixel value and

\begin{equation}
\mathrm{MSE} = \frac{1}{N} \sum_{i=1}^{N} 
(\hat{S}_i - S_i)^2.
\end{equation}

Higher PSNR values indicate better fidelity. Although PSNR does not explicitly model perceptual similarity, it remains a standard metric for quantitative comparison in image restoration studies.

\subsection{Structural and Perceptual Metrics}


The \textbf{Structural Similarity Index }(SSIM) is a perceptual metric designed to measure similarity between two images based on luminance, contrast, and structural information. Unlike pixel-wise losses, SSIM evaluates local structural consistency by comparing statistics computed over small neighborhoods.

For two image patches $x$ and $y$, SSIM is defined as

\begin{equation}
\mathrm{SSIM}(x, y) = 
\frac{(2\mu_x \mu_y + C_1)(2\sigma_{xy} + C_2)}
{(\mu_x^2 + \mu_y^2 + C_1)(\sigma_x^2 + \sigma_y^2 + C_2)},
\end{equation}

where $\mu_x$ and $\mu_y$ denote local means, $\sigma_x^2$ and $\sigma_y^2$ denote variances, $\sigma_{xy}$ denotes covariance, and $C_1$, $C_2$ are small constants for numerical stability. SSIM values range between $-1$ and $1$, with higher values indicating greater structural similarity.

SSIM is particularly suitable for medical imaging applications, where preservation of structural information is critical.

\textbf{Visual Information Fidelity (VIF)} is an information-theoretic perceptual metric that quantifies the amount of visual information preserved in the reconstructed image relative to the reference. It is defined as:
\begin{equation}
\mathrm{VIF} =
\frac{\sum_{i} I(C_i; F_i)}
{\sum_{i} I(C_i; E_i)}
\end{equation}
where $I(\cdot;\cdot)$ denotes mutual information, $C_i$ represents the reference image information, $F_i$ the distorted image, and $E_i$ the perceived image through the human visual system model.

\textbf{Learned Perceptual Image Patch Similarity (LPIPS)} is a deep feature-based perceptual metric that measures similarity between images using features extracted from pretrained neural networks. It is computed as:
\begin{equation}
\mathrm{LPIPS}(x, y) =
\sum_{l} \frac{1}{H_l W_l}
\sum_{h,w} 
\left( w_l \cdot \left( \phi_l(x)_{h,w} - \phi_l(y)_{h,w} \right) \right)^2
\end{equation}
where $\phi_l$ represents deep features from layer $l$, and $w_l$ are learned weights. Lower LPIPS values indicate better perceptual similarity.
 
 


\subsection{Computational Efficiency Metrics}

In addition to reconstruction quality, the efficiency of a model is an important factor, particularly for high-resolution CT data and potential clinical deployment. 

The number of Parameters (\textbf{Params}) is the total count of trainable weights in the network and reflects the model’s memory footprint as well as its representational capacity.

\textbf{Floating Point Operations (FLOPs)} quantify the number of arithmetic operations required to perform a forward pass through the network. FLOPs are computed from multiply–accumulate operations (MACs) obtained using the ptflops library, where each MAC is counted as two floating point operations (one multiplication and one addition). This metric provides an estimate of the computational complexity of the model.

Furthermore, \textbf{Inference time }measures the actual runtime required to process an input sample, typically reported in milliseconds. Unlike FLOPs, inference time reflects real-world performance and depends on both the model architecture and hardware implementation.




In practice, no single metric captures all aspects of image quality. Pixel-wise measures like L1 and PSNR focus on numerical accuracy between restored and reference images, whereas perceptual metrics like SSIM better reflect structural similarities that matter most for clinical image interpretation. For this reason, image restoration studies typically report several complementary metrics together. 
\begin{quote}


\end{quote}

 

\begin{quote}


\end{quote}

\begin{quote}


\end{quote}

 

\begin{quote}


\end{quote}

\begin{quote}


\end{quote}

 


   % (\chapter{})
\cleardoublepage
\chapter{Related Work}
\label{chap:relatedwork}
\section{Image-Domain Restoration Methods}

Deep learning methods have been widely applied to CT image restoration, leveraging convolutional neural networks to learn complex mappings between degraded and clean images through hierarchical feature extraction.

Among the earliest CNN-based approaches for CT, Chen et al. \cite{Che+17a} proposed RED-CNN, a residual encoder–decoder network for low-dose CT denoising that demonstrated the effectiveness of combining residual learning with encoder–decoder architectures for medical image restoration. 

More recently, frequency-aware designs have gained attention in image restoration. Chen et al. \cite{LKFN2024} proposed a large-kernel frequency-enhanced network that uses spectral analysis to emphasize high-frequency components corresponding to edges and fine structural details. Hu et al. \cite{hu2025ifsr} introduced an implicit frequency selection mechanism that enables networks to adaptively prioritize informative frequency bands during restoration. 

\section{Projection-Domain Restoration} 
Projection-domain approaches aim to correct artifacts directly in sinogram space before image reconstruction. Since artifacts originate from inconsistencies in projection measurements, addressing them at this stage can prevent error propagation into reconstructed images.

Existing work in this area has primarily focused on metal artifact reduction. Ghani and Karl \cite{Gha2018} applied a deep learning-based network directly in the sinogram domain, demonstrating that operating on the entire sinogram rather than local patches is essential due to the non-local nature of sinogram information. Podgorsak et al. \cite{Pod2021} applied DCGANs for sinogram-domain correction of sparse and truncated projections, showing that sinogram-domain correction outperformed image-domain correction for sparse acquisitions.




\section{Model-Based and Optimization-Based CT Reconstruction}
Traditional approaches for artifact reduction in computed tomography rely on model-based and optimization-driven frameworks. These methods define motion correction as an inverse problem, where the goal is to estimate motion parameters and reconstruct images that are consistent with the measured projection data \cite{ritschl2011motion}.
A common strategy involves a two-step pipeline of motion estimation followed by motion-compensated reconstruction. Hahn et al. \cite{Hah+17}  developed a motion compensation framework for cone-beam CT that estimates rigid motion parameters from the projection data and incorporates them into the backprojection operator. Chee et al. \cite{Che2019} proposed MC-SART, which iteratively refines both the motion model and the reconstructed volume directly from cone-beam projections. More recently, Preuhs et al. \cite{Pre+20} introduced a differentiable approach to motion parameter estimation in flat-panel CT, enabling gradient-based optimization of rigid motion parameters within the reconstruction pipeline. Building on this direction, Thies et al. \cite{thies2025_cbctmotion} developed a fully differentiable imaging model that computes analytical gradients of the reconstruction operator for efficient head motion compensation in cone-beam CT.

While these methods are grounded in physical modeling and provide accurate results, they are typically computationally expensive and depend on accurate motion estimation, which may be difficult to achieve in practical clinical scenarios.


\section{Unsupervised and Image-to-Image Translation Methods  }

In many clinical scenarios, paired datasets of artifact-corrupted and clean CT images are difficult to obtain, motivating the development of unsupervised and semi-supervised restoration approaches \cite{zhu2017cyclegan}.  

In the context of CT artifact reduction, Liao et al. \cite{liao2019adn} proposed an artifact disentanglement network that separates artifact components from underlying anatomical structures in a learned latent space, enabling unsupervised metal artifact removal without paired training data.

Although unsupervised methods improve generalization and reduce dependency on synthetic data, they are typically applied in the image domain and may not fully address inconsistencies originating in projection measurements.



 


\section{Transformer-Based Methods} 

Transformer-based architectures have been introduced to address the limitations of CNNs in modeling long-range dependencies. By employing self-attention mechanisms, transformers can capture global context across the entire image.

For image restoration, several transformer-based models have been developed. Liang et al. \cite{liang2021swinir} proposed SwinIR, which applies window-based self-attention with shifted windows to enable efficient cross-window interaction while maintaining linear computational scaling with respect to image size. Zamir et al. \cite{restormer2022} introduced Restormer, which computes self-attention across channels rather than spatial dimensions, enabling efficient processing of high-resolution images. Wang et al. \cite{Wan+22} proposed Uformer, combining a U-shaped hierarchical structure with window-based self-attention for efficient multi-scale image restoration. 

Despite their strong modeling capabilities, transformer-based architectures incur higher computational complexity and memory consumption compared to convolutional approaches. This limitation motivates the exploration of alternative architectures that can capture long-range dependencies more efficiently.



\section{Receptive Field Expansion and Large-Kernel CNNs}

A key factor in modeling global dependencies is the size of the receptive field. Traditional CNNs rely on stacking multiple small convolutional kernels to gradually increase the receptive field \cite{Luo2016}. However, this approach is inefficient and may not effectively capture long-range dependencies.

Recent works revisit large-kernel convolution as a more direct approach. Ding et al. \cite{replknet2022} proposed RepLKNet, demonstrating that depthwise convolution kernels as large as 31×31 can be employed efficiently and achieve performance comparable to transformer-based models on various vision tasks. Liu et al. \cite{Liu2022} showed with ConvNeXt that systematically modernizing a standard ResNet with transformer-inspired design choices, including depthwise convolutions with 7×7 kernels, inverted bottlenecks, and layer normalization, enables pure convolutional architectures to match or exceed Swin Transformer performance. Building on this, Liu et al. \cite{Liu2023} proposed SLaK (Sparse Large-Kernel Network), which uses dynamic sparsity and kernel decomposition to scale kernels to 51×51, overcoming the performance saturation observed in RepLKNet beyond 31×31. More recently, Ding et al.  \cite{Din+24} introduced UniRepLKNet, a universal architecture that further scales kernels and demonstrates the effectiveness of large-kernel CNNs across diverse tasks, including image classification, semantic segmentation, and object detection. 

To address the computational challenges associated with large-kernels, several efficient implementations have been proposed. Asymmetric kernel decomposition approximates large square kernels using smaller rectangular convolutions, reducing computational cost while preserving spatial coverage \cite{qu2024alkdnet}. Chen et al. \cite{PeLK2024} proposed parameter-efficient large-kernel convolutions using peripheral convolution, which captures large spatial context with substantially fewer parameters than standard large-kernel depthwise convolutions.


These approaches highlight that expanding the receptive field is a fundamental requirement for improving performance in image restoration tasks.

\section{Multi-scale and Hybrid Kernel Designs}
Recent research has explored combining multiple kernel sizes and feature representations within a unified framework. Multi-scale architectures integrate local, large, and global receptive fields to capture features at different spatial scales \cite{ronneberger2015unet}.

Omni-kernel networks introduce parallel branches that process global, large-scale, and local features simultaneously, enabling the model to handle diverse types of degradations \cite{cui2024oknet}. Additionally, hybrid approaches combine spatial and frequency-domain representations to improve feature extraction and robustness \cite{LKFN2024}.

These designs demonstrate that combining multiple receptive field scales is beneficial for capturing both fine details and global structures in image restoration tasks.

\section{Summary and Research Gap}
In summary, existing methods for artifact reduction and image restoration can be categorized into image-domain, projection-domain, and model-based approaches. Image-domain methods are effective but limited by their inability to address inconsistencies at the source. Model-based approaches are physically grounded but computationally expensive. Projection-domain methods offer a promising direction but require models capable of capturing global dependencies.

Recent developments, including transformer-based models and large-kernel convolutional networks, have shown strong results in capturing long-range dependencies, which is important for dealing with motion artifacts.  

However, the scope of super-large-kernel convolutional networks for motion-artifact correction remains largely unexplored. Whether such architectures can effectively combine the computational efficiency of convolutional designs with the global modeling capacity required for artifact correction remains an open question. 
 


 % (\chapter{})
\cleardoublepage
\chapter{Models}
\label{chap:models}
\section{Overview of Considered Models}
This work investigates the effectiveness of super-large-kernel convolution for motion artifact reduction in CT data. 
For effective comparison, a combination of convolutional, large-kernel, and transformer-based architectures is considered to evaluate performance across different receptive fields and modeling capabilities.

The models included are:

\begin{itemize}
    \item Proposed model: \textbf{MR-LKV (Motion Restoration using Large Kernel conVolution)} 
        \item Baselines: 
    \begin{itemize}
        \item U-Net ~\cite{ronneberger2015unet}
        \item RepLKNet ~\cite{replknet2022}
        \item SwinIR  ~\cite{liang2021swinir}
        \item Restormer ~\cite{restormer2022}
    \end{itemize}
\end{itemize}
\subsection{Domain-wise Model Application}

Although identical baseline architectures are used across both domains, the proposed model is modified according to the domain in which the training is employed. 
The domain-wise model assignments are summarized in Table \ref{tab:domain_comparison}. 

\begin{table}[htpb]
\centering
\caption{Comparison of Projection-Domain and Image-Domain Training Setups}
\small
\begin{tabular}{|c|c|c|p{6cm}|}
\hline
\textbf{Domain} & \textbf{Input} & \textbf{Output} & \textbf{Models} \\ \hline

Projection & 2D sinogram & Clean sinogram & MR-LKV (full), U-Net, RepLKNet, SwinIR, Restormer \\ \hline

Image & CT slices & Clean CT & MR-LKV (simplified), same baselines \\ \hline

\end{tabular}
\label{tab:domain_comparison}
\end{table}
 \section{Proposed Model: MR-LKV}
 \subsection{Design Motivation}
Motion artifacts in CT imaging exhibit complex spatial patterns that span large regions and often arise from inconsistencies across projection views. Conventional convolutional neural networks with small kernel sizes struggle to model such global structures efficiently ~\cite{Luo2016}.  Transformers effectively capture long-range dependencies but incur computational costs that limit their practicality in real-time clinical systems \cite{restormer2022}.

To address these limitations, the proposed MR-LKV model is designed with three main objectives:

\begin{itemize}
    
    \item Capture \textbf{long-range spatial dependencies} using super-large-kernel convolutions 
    \item Incorporate \textbf{frequency-domain processing} to correct global inconsistencies 
    \item Integrate \textbf{attention mechanisms} to emphasize informative spatial and directional features


\end{itemize}

The architecture has a flexible framework that can be adapted for both projection-domain and image-domain restoration tasks. 
\subsection{Architectural Overview}

The proposed MR-LKV architecture is a multi-scale encoder--decoder network designed to capture both global and local dependencies in degraded CT data. The model integrates frequency-domain processing, attention mechanisms, and super-large-kernel convolutional blocks within a unified framework.

The architecture follows a U-Net-like structure ~\cite{ronneberger2015unet} consisting of:

\begin{itemize}
    \item an initial patch embedding layer,
    \item four encoder stages with hierarchical feature extraction,
    \item a symmetrical decoder with skip connections, and
    \item a final reconstruction layer with residual learning.
\end{itemize}

The complete architecture is shown in figure \ref{fig:mrlkv_architecture}

Each encoder stage applies a sequential feature refinement pipeline before spatial encoding, ensuring that informative features are emphasized prior to super-large-kernel processing.

\subsubsection{Patch Embedding and Multi-Scale Representation}

Given an input $S \in \mathbb{R}^{1 \times H \times W}$, an initial convolutional embedding layer is applied:

\begin{equation}
X_0 = \mathrm{Conv}_{3\times3,\, s=2}(S)
\end{equation}

This operation reduces spatial resolution while increasing the number of feature channels. 
The network then follows a hierarchical structure with channel dimensions:

\begin{equation}
C = \{32, 64, 128, 256\}
\end{equation}

At each stage, spatial resolution decreases while channel capacity increases, enabling the network to capture both fine-grained and global features.

\subsubsection{Encoder Design and Feature Refinement Pipeline}

Each encoder stage consists of a structured sequence of operations:

\begin{equation}
X \rightarrow \mathrm{FEB} \rightarrow \mathrm{CA} \rightarrow \mathrm{View\text{-}Attn} \rightarrow \mathrm{RLKB}
\end{equation}

followed by downsampling (except for the final stage).

This sequential design ensures that global contextual information and attention-based feature selection are incorporated prior to spatial modeling.

\subsubsection{Frequency-Enhanced Branch (FEB)}

At each encoder stage, a frequency-enhanced branch (FEB) is applied in the beginning to capture global inconsistencies that are difficult to model using purely spatial operations.

Given an input feature map $X \in \mathbb{R}^{C \times H \times W}$, the FEB first applies a $1 \times 1$ convolution for channel-wise feature transformation:

\begin{equation}
X' = \mathrm{Conv}_{1\times1}(X)
\end{equation}


This transformation enables channel mixing before frequency-domain processing. The transformed features are then processed using a 2D FFT:

\begin{equation}
\hat{X} = \mathrm{FFT}_{2\mathrm{D}}(X')
\end{equation}

Separate $1 \times 1$ convolutions are applied to the real and imaginary components.


The transformed features are processed in the frequency domain and mapped back:


\begin{equation}
    Y = \mathrm{IFFT}_{2\mathrm{D}}(f(\hat{X}))
\end{equation}

A second $1 \times 1$ convolution, Batch Normalization, and residual connection with GELU activation produce:

\begin{equation}
X_{\mathrm{feb}} =
\mathrm{GELU}\left(
\mathrm{Norm}(\mathrm{Conv}_{1\times1}(Y)) + X
\right)
\end{equation}

This module enables the network to model structured distortions in the frequency domain.

\subsubsection{Attention Mechanisms}

Two complementary attention mechanisms are applied sequentially.

\paragraph{Channel Attention (CA)}
Channel attention emphasizes informative feature channels using global context 

\begin{equation}
X_{\mathrm{ca}} = X \odot w
\end{equation}

where $w$  is obtained from global average pooling followed by a lightweight two-layer network with reduction ratio $r = 16$.


\paragraph{View-Axis Attention}

To model directional dependencies, attention is applied along the width dimension $W$, which corresponds to projection views.

\begin{equation}
X_{\mathrm{vaa}} = X \odot a
\end{equation}
where $a$ is obtained via pooling along the detector axis followed by 1D convolution along the view axis.

The intermediate channel dimension is reduced by a factor of 8.

\subsubsection{Residual Large-Kernel Blocks (RLKB)}
The RLKB is the core component of MR-LKV, employing progressively increasing kernel sizes to capture long-range dependencies  
The RLKB performs spatial modeling using large receptive fields.

\begin{equation}
U = \mathrm{PW}(\mathrm{DW}(X))
\end{equation}

\begin{equation}
Y =
\mathrm{GELU}\left(
\mathrm{Norm}(U) + X
\right)
\end{equation}

where DW denotes depthwise convolution and PW denotes pointwise convolution.

\paragraph{Implementation Details}
\begin{itemize}
    \item Kernel sizes: $35, 55, 75, 95$ (stage-wise)
    \item Channels: $32, 64, 128, 256$
    \item Depthwise convolution: groups = $C$
    \item Padding preserves spatial dimensions
    \item Normalization: BatchNorm2d
    \item Activation: GELU
\end{itemize}

\paragraph{Design Rationale}

Large kernels enable direct modeling of long-range dependencies. Depthwise separable design ensures computational efficiency while maintaining large receptive fields.

\subsubsection{Downsampling}

Between encoder stages, downsampling is performed using:

\begin{equation}
X_{i+1} =
\mathrm{Conv}_{2\times2,\, s=2}(X_i)
\end{equation}

This reduces spatial resolution while doubling channel dimensions.

\subsubsection{Decoder}

The decoder reconstructs the output using bilinear upsampling and skip connections.

\begin{equation}
Y =
\mathrm{Conv}_{1\times1}\left(
[\mathcal{U}(X), S]
\right)
\end{equation}

where $\mathcal{U}(X)$ is upsampled features and $S$ is the corresponding encoder feature map.

Skip connections preserve spatial detail, and the decoder avoids large-kernel operations to reduce computational cost.

\subsubsection{Output Formulation}

The final output is obtained using residual learning:

\begin{equation}
\hat{S} = S + f_\theta(S)
\end{equation}

where $f_\theta$ represents the MR-LKV network.

This formulation allows the model to focus on learning artifact corrections rather than reconstructing the entire signal.


The detailed structure of each component is illustrated in Fig. \ref{fig:depth_mrlkv_architecture} 

\subsection{Domain-Specific Adaptation of MR-LKV }

While the MR-LKV architecture is designed as a unified framework for modeling global motion artifacts, its implementation is adapted to suit the characteristics of projection-domain and image-domain data. For the projection-domain, the architecture remains unchanged, whereas for the image-domain, there has been some modifications.

\subsubsection{Removal of Frequency-enhanced Branch}
    The \textbf{frequency-enhanced branch} is omitted in the image-domain model.
In reconstructed CT images:

\begin{itemize}
    \item artifacts appear as spatial distortions rather than structured frequency inconsistencies 
    \item global frequency-domain modeling provides limited additional benefit 
\end{itemize}
Thus, the feature refinement pipeline is simplified by excluding the FEB.
\subsubsection{Removal of View-Axis Attention}
The \textbf{view-axis attention} mechanism is also removed in the image-domain configuration.

Unlike sinograms, reconstructed CT images:

\begin{itemize}
    \item do not contain an explicit directional axis corresponding to projection views 
    \item exhibit isotropic spatial structure 
\end{itemize}
As a result, applying directional attention along a specific axis is not meaningful and may introduce unnecessary bias.

\subsubsection{Retained Components}
Despite these modifications, the core architectural principles of MR-LKV are preserved.

The image-domain model retains:
\begin{itemize}
    \item  Residual Large-Kernel Blocks (RLKB)
    \item  Channel Attention (CA)
    \item Encoder--decoder structure with skip connections
    \item Multi-scale feature representation
    \item Residual learning at output
\end{itemize}




The processing pipeline in the image-domain becomes: 

\begin{equation}
X \rightarrow  \mathrm{CA} \rightarrow \mathrm{RLKB}
\end{equation}
This simplified structure focuses on spatial feature modeling while maintaining efficiency.

The full configuration of the MR-LKV architecture is provided in Table \ref{tab:mrlkv_architecture} 

\begin{table}[H]
\centering
\caption{Configuration of the proposed MR-LKV architecture}
\label{tab:mrlkv_architecture}
\renewcommand{\arraystretch}{1.2}
\begin{tabular}{l c c c p{5.5cm}}
\toprule
\rowcolor{gray!15}
\textbf{Stage} & \textbf{Channels} & \textbf{RLKB Blocks} & \textbf{Kernel Size} & \textbf{Operations} \\
\midrule
Patch Embed & 32 & -- & $3 \times 3$ & Stride-2 convolution (downsampling)\\
Stage 0 & 32 & 2 & $35 \times 35$ & (FEB → CA → View-Attn) + RLKB\\
Stage 1 & 64 & 2 & $55 \times 55$ & (FEB → CA → View-Attn) + RLKB\\
Stage 2 & 128 & 3 & $75 \times 75$ & (FEB → CA → View-Attn) + RLKB\\
\rowcolor{green!10}
\textbf{Stage 3} & \textbf{256} & \textbf{2} & $\mathbf{95 \times 95}$ & (FEB → CA → View-Attn) + RLKB\\
\midrule
Decoder (Stage 2) & 128 & -- & $1 \times 1$ & Bilinear upsampling + skip concatenation + 1×1 convolution\\
Decoder (Stage 1) & 64 & -- & $1 \times 1$ & Bilinear upsampling + skip concatenation + 1×1 convolution\\
Decoder (Stage 0) & 32 & -- & $1 \times 1$ & Bilinear upsampling + skip concatenation + 1×1 convolution\\
Output & 1 & -- & $1 \times 1$ & Final projection + residual addition\\
\bottomrule
\end{tabular}
\end{table}
 \begin{figure}[htpb]
    \centering
    \includegraphics[width=1.0
    \linewidth]{MRLKV.drawio.png}
    \caption{Overall architecture of the proposed MR-LKV network. The encoder consists of multiple stages, each composed of a Frequency Enhancement Branch (FEB), channel attention, view-axis attention, and stacked Residual Large-Kernel Blocks (RLKBs) with progressively increasing kernel sizes to capture long-range dependencies. }
    \label{fig:mrlkv_architecture}

\end{figure}

\begin{figure}[htpb]
    \centering
    \includegraphics[width=1.0\linewidth]{building blocks.pdf}\caption{Key components of the MR-LKV architecture: frequency enhancement, sequential attention, and residual large-kernel convolution blocks. }
   
    \label{fig:depth_mrlkv_architecture}
\end{figure}






\subsection{Summary}

The MR-LKV architecture is designed as a flexible framework that can be adapted based on the domain. By removing domain-specific components such as frequency enhancement and view-axis attention, the image-domain model achieves a more efficient and targeted design for spatial artifact correction while retaining the benefits of large-kernel convolutional modeling. 

\section{Baseline Architecture: U-Net}

U-Net is a convolutional encoder--decoder architecture with skip connections, widely used for image restoration tasks ~\cite{ronneberger2015unet}. It is included as a baseline to represent conventional small-kernel CNN approaches for both projection-domain and image-domain artifact correction.

\subsubsection{Architecture}
A modified implementation of the standard U-Net ~\cite{ronneberger2015unet} is used with additional engineering improvements for stability.
The network follows a symmetric encoder--decoder structure with four resolution levels. Each encoder stage consists of two convolutional layers with $3 \times 3$ kernels, followed by downsampling using max-pooling. The feature dimensionality grows progressively across the network stages, allowing hierarchical extraction of representations at different levels of abstraction.


At the lowest resolution, a bottleneck block further processes the features before passing them to the decoder. The decoder reconstructs spatial details using transposed convolutions, while skip connections are employed to fuse encoder and decoder features, preserving spatial information across scales.

\subsubsection{Implementation Details}

In the conventional U-Net configuration, the base number of feature channels is typically set to 64. However, in this work, the architecture is modified to use a reduced base width:

\begin{itemize}
    \item Input/output channels: 1
    \item Base channels: 32 (reduced from standard 64)
    \item Number of levels: 4
    \item Convolution kernel size: $3 \times 3$
    \item Normalization: BatchNorm
    \item Bottleneck dropout: 0.1
    \item Weight initialization: Xavier uniform initialization for convolution layers 
    \item  Skip connection alignment: Center cropping is used to handle spatial mismatches 
\end{itemize}

The reduction in base channels decreases the overall number of parameters and computational cost, enabling a fair comparison with other models under similar resource constraints.

\subsubsection{Discussion}

Owing to its use of small convolutional kernels, U-Net captures local structures effectively but requires deeper feature hierarchies to model long-range dependencies. This constraint is particularly relevant in image space, where motion artifacts exhibit globally distributed patterns.

\section{Baseline Architecture: RepLKNet}

RepLKNet is a convolutional architecture designed to capture long-range spatial dependencies using large-kernel depthwise convolutions ~\cite{replknet2022}. It is included as a baseline to evaluate the effectiveness of large-kernel CNNs in removing CT artifacts.

\subsubsection{Architecture}

The configuration follows an encoder–decoder architecture, organized into multiple hierarchical stages. Each stage contains several large-kernel convolutional blocks in which depthwise convolutions with expanded receptive fields are followed by pointwise convolutions for inter-channel interaction.  Skip connections are incorporated between corresponding encoder and decoder layers to preserve spatial information. 

To improve training stability, each block integrates a parallel small-kernel depthwise branch that operates alongside the large-kernel convolution, and the two outputs are combined within a residual connection.  Additionally, convolutional feed-forward (ConvFFN) layers are interleaved between RepLK blocks to enhance feature representation.

\subsubsection{Implementation Details}

RepLKNet is adapted for image restoration with an encoder--decoder structure and the following configuration:

\begin{itemize}
    \item Input/output channels: $1$
    \item Stage channels: $[32, 64, 128, 256]$
    \item Number of blocks per stage: $[2, 2, 3, 2]$
    \item Large kernel sizes: $[21, 19, 17, 9]$
    \item Small kernel size: $5$
    \item Normalization: $GroupNorm$ (1 group), enabling stable training for small batch sizes 
    \item  Feature refinement: ConvFFN blocks alternated with RepLK blocks  
    \item Residual learning: Applied at both block and stage levels  
    \item  Multi-scale kernel design: Large and small depthwise convolutions are combined  
    \item Feature alignment: Bilinear interpolation is used to match feature map sizes during decoding
\end{itemize}

Compared to the original RepLKNet design ~\cite{replknet2022}, the model is modified in several aspects to make it suitable for image restoration task as compared to classification task which was the intended nature of the original format. 

To ensure compatibility with the hierarchical structure, the input resolution is constrained to be divisible by 4, and inputs are cropped accordingly during training.

\subsubsection{Discussion}

RepLKNet directly models long-range dependencies through large convolutional kernels, making it advantageous for capturing globally distributed motion artifacts. MR-LKV improves upon this by incorporating an increasing super-large-kernel setup, leading to better feature refinement and artifact suppression. 


 \section{Baseline Architecture: SwinIR}

SwinIR is a transformer-based framework for image restoration that applies windowed self-attention to learn local dependencies within fixed regions while approximating global context through shifted window interactions across layers ~\cite{liang2021swinir}. It is included as a benchmark to evaluate attention-based approaches for artifact correction.
\subsubsection{Architecture}

The model consists of stacked Residual Swin Transformer Blocks (RSTBs), where self-attention is computed within local windows and shifted across layers to enable cross-window interaction. Each RSTB contains several Swin Transformer blocks followed by a convolutional layer within a residual connection. 

The overall structure begins with a shallow feature extraction layer, proceeds through deeper transformer-based representation learning, and concludes with a reconstruction module. Residual learning is integrated throughout the design to stabilize optimization and to preserve low-frequency image components.

\subsubsection{Implementation Details}

In this work, SwinIR is configured for CT restoration with the following settings:

\begin{itemize}
    \item Input/output channels: $1$
    \item Embedding dimension: $48$
    \item Number of transformer blocks per RSTB: $[2, 2, 2, 2]$
    \item Number of attention heads: $[2, 2, 2, 2]$
    \item Window size: $8$
    \item Patch size : $128 \times 128$
    \item Batch size : 2
    \item Upsampling: $disabled$
\end{itemize}

Relative to the original SwinIR design ~\cite{liang2021swinir}, both the embedding dimension and depth are reduced to improve computational efficiency for high-resolution inputs. Additionally, checkpointing is enabled during training to reduce memory usage.

To ensure compatibility with the hierarchical structure, the input resolution is constrained to be divisible by 8, and inputs are cropped accordingly during training.

\subsubsection{Discussion}

SwinIR captures both local and global dependencies through window-based attention and shifted window interactions. Nonetheless, the use of self-attention increases computational complexity, compared to convolutional approaches.


\section{Baseline Architecture: Restormer}

Restormer is a transformer-based architecture developed for high-resolution image restoration, combining self-attention with convolution-based operations embedded within transformer blocks ~\cite{restormer2022}. It is included as a baseline to evaluate multi-scale attention-based models for artifact correction.

\subsubsection{Architecture}

The framework follows a hierarchical encoder--decoder structure operating across multiple resolution levels. Each level consists of transformer blocks that fuse multi-head self-attention with convolutional feed-forward networks ( \textbf{Gated-DConv Feed-Forward Network)}. 
A key innovation of Restormer is the \textbf{Multi-DConv Head Transposed Attention (MDTA)} module, where attention computation employs depthwise and pointwise convolutions. This structure enhances local feature extraction while maintaining efficiency. 
Downsampling and upsampling operations are used to capture multi-scale contextual information, while skip connections are employed to preserve spatial features across stages.
Unlike standard transformers, Restormer embeds depthwise and pointwise convolutions directly within both the attention and feed-forward modules to improve computational efficiency and local feature representation 

\subsubsection{Implementation Details}

In this work, Restormer is adapted for CT image restoration with the following configuration:

\begin{itemize}
    \item Input/output channels: 1
    \item Base feature dimension: 48
    \item Number of blocks per level: [2, 2, 2, 3]
    \item Number of attention heads: [1, 2, 4, 8]
    \item Feed-forward expansion factor: 2.0
    \item Number of refinement blocks: 2
    \item Patch Size: $128 \times 128$
    \item Batch Size: 1
\end{itemize}

Compared to the original implementation ~\cite{restormer2022}, the model is simplified by reducing the number of blocks per stage and the feed-forward expansion factor to improve computational efficiency. Additionally, the network is adapted for both single-channel sinogram inputs and CT image slice inputs.

To ensure compatibility with the hierarchical structure, the input resolution is constrained to be divisible by 8, and inputs are cropped accordingly during training.

\subsubsection{Discussion}

Restormer captures long-range dependencies through attention mechanisms while leveraging convolutional operations for efficient feature representation. Its multi-scale design enables effective modeling of structured artifacts, but it introduces higher computational and memory requirements compared to convolutional architectures.
\cleardoublepage
\chapter{Methodology}
\label{chap:methodology}
\section{ Overview of the Proposed Approach} 


The goal of this work is to reduce motion artifacts in reconstructed CT images using a dual pipeline framework (image-domain and projection-domain). The overall methodology consists of three common stages followed by the two domain-specific restoration pipelines. The backbone is a CNN that incorporates super‑large-kernels convolution so that individual layers can capture the long‑range spatial patterns that are characteristic of motion artifacts. 

First, clean 3D CT volumes are converted into projection data (sinograms) using forward projection \cite{diffct2025}. Motion artifacts are then introduced directly in the projection-domain to simulate realistic acquisition inconsistencies. These corrupted sinograms serve as the basis for both restoration pipelines.
Following artifact simulation, two complementary restoration strategies are explored.  
In the \textbf{projection-domain} pipeline, models are trained to directly correct motion-corrupted sinograms before reconstruction (after converting them to 2D sinograms), whereas, in the \textbf{image-domain }pipeline, both clean and corrupted sinograms are first reconstructed into CT images using cone-beam filtered backprojection \cite{diffct2025}, and restoration is performed directly in the image-domain.

\paragraph{Problem Formulation}
Let $y$ denote a clean reference and $x $ its motion-corrupted counterpart. 
The objective of the proposed method is to learn a parametric mapping:

\begin{equation}
\hat{y} = f_\theta(x)
\end{equation}

where $f_\theta$ represents the MR-LKV network with parameters $\theta$, and $\hat{y}$ denotes the restored output. 
In this work, the formulation is applied in both projection- and image-domains, enabling the model to learn either sinogram-to-sinogram or image-to-image mappings depending on the pipeline.


An overview of the proposed framework is illustrated in Fig. \ref{fig:pipeline}, highlighting the shared preprocessing stages and the divergence into projection-domain and image-domain learning frameworks. 

\begin{figure}[htpb]
    \centering
    \includegraphics[width=1.0\linewidth]{pipeline.pdf}
    \caption{Pipeline of the proposed MR-LKV framework showing 3D sinogram generation, motion artifact simulation, domain-specific 2D training and inference, reconstruction, and volumetric evaluation. }
    \label{fig:pipeline}
\end{figure}


 \section{Dataset and Preprocessing}
 \subsection{Dataset Description}
The experiments in this thesis are conducted using CT data from the \textbf{CQ500} dataset \cite{CQ500}, consisting of \textbf{491} non-contrast head CT scans acquired under clinical conditions.

From this dataset, an initial subset of 100 patients is selected. For each patient, the relevant CT series is identified by scanning all subdirectories within the patient folder and selecting only those whose names match a predefined set of valid series descriptors: specifically, \texttt{ct plain}, \texttt{ct plain thin}, or \texttt{ct thin plain} (after normalizing underscores, hyphens, and whitespace to lowercase). 
Any series containing fewer than 20 DICOM files is excluded. When multiple valid series exist for a single patient, the one with the largest number of slices is selected. After applying this criteria, \textbf{53 patients} are retained for further processing, with one CT series selected per patient.
These 53 patients are used exclusively for model training and validation. A separate, non-overlapping set of 30 patients from the remaining CQ500 data is set for final evaluation. 


For each selected series, Digital Imaging and Communications in Medicine (DICOM) slices are loaded and sorted according to their \textbf{InstanceNumber} to ensure correct anatomical ordering. The slices are then stacked into 3D volumetric arrays of shape $(Z, Y, X)$, where $Z$ denotes the number of slices (axial dimension) and $(Y, X)$ represents the in-plane spatial dimensions (height and width).

\subsection{DICOM Volume Normalization}
The CT volumes are loaded directly from DICOM files using raw pixel intensity values without applying rescale slope and intercept to convert them into Hounsfield Units (HU). Since the proposed method operates on normalized data and focuses on learning relative intensity patterns for artifact reduction, absolute HU calibration is not required. However, this assumption may affect generalizability, as raw pixel intensities can vary across scanners and acquisition protocols. 
To maintain numerical stability and consistent scaling across samples, each 3D volume is independently normalized in the range $[0,1]$ using min–max normalization defined as:

\begin{equation}
x_{\mathrm{norm}} = \frac{x - x_{\min}}{x_{\max} - x_{\min}}
\end{equation}
where $x_{min}$ and $x_{max}$ denote the minimum and maximum voxel intensities within each volume.

\subsection{Conversion from DICOM to Sinogram}

Since motion corruptions occur in the projection-domain, reconstructed CT images in DICOM format are first converted into sinogram data. For this purpose, a cone‑beam forward projection operator from the Differentiable CT (\textbf{DiffCT}) library is used \cite{diffct2025} on a 3D cone-beam configuration, producing volumetric sinograms that preserve spatial consistency across slices. 
    
In this work, the DiffCT implementation is used as a preprocessing tool without modifying its internal representation. However, key geometric and detector-related parameters, such as detector pixel spacing ($du, dv$), source-to-isocenter distance (SID), and source-to-detector distance (SDD), are adjusted to match the spatial resolution and field of view of the input CT images. The voxel spacing for the forward projection is extracted from the DICOM metadata. The in-plane pixel spacing is read from the \texttt{PixelSpacing} attribute, and the slice thickness is obtained from the \texttt{SliceThickness} attribute (defaulting to the minimum of the in-plane spacings if unavailable). To simplify the projection geometry, a single scalar voxel spacing is defined as the minimum of these values and used uniformly during forward projection. The projection geometry parameters are listed in Table \ref{tab:projection_geometry}. 

The projection geometry includes a fixed number of 540 projection views uniformly distributed over a full angular range. Each projection is computed on GPU using the differentiable cone-beam forward projector as described above, producing 3D sinograms of shape $(N\_{views}, U, V)$, where $U$ and $V$denote the horizontal and vertical detector dimensions, respectively. For each patient, the resulting clean 3D sinogram and a corresponding visualization are saved for subsequent processing. 


\begin{table}[htbp]
\centering
\caption{Cone-beam projection geometry parameters used for sinogram generation.}
\label{tab:projection_geometry}
\begin{tabular}{lc}
\hline
\textbf{Parameter} & \textbf{Value} \\
\hline
Detector resolution & $800 \times 800$ \\
Detector pixel spacing ($du, dv$) & $1.0$ \\
Source-to-isocenter distance (SID) & $530$ \\
Source-to-detector distance (SDD) & $1095$ \\
Number of projection views & $540$ \\
Angular coverage & $360^\circ$ \\
\hline
\end{tabular}
\end{table}

This preprocessing pipeline is summarized in Fig. \ref{fig:dicom_to_sinogram}. 

\begin{figure}[htpb]
    \centering
   
        \includegraphics[width=1.0\linewidth]{dicom_to_sino.pdf}
       \caption{Overview of the preprocessing pipeline from DICOM volume to 3D Sinogram.}
        \label{fig:dicom_to_sinogram}
    \end{figure}
   

\section{Motion Artifact Simulation in 3D Sinogram}

Motion artifacts in CT occur when projections acquired at different time points no longer correspond to the same object position during image acquisition \cite{wagner2003quantification}. In this work, motion artifacts are simulated by applying geometric transformations directly to the clean 3D sinograms produced in Section~5.2.3, generating a corresponding set of motion-corrupted 3D sinograms \cite{thies2025_cbctmotion}. 
Let
\begin{equation}
S \in \mathbb{R}^{N\_{views} \times U \times V}
\end{equation}
denote an artifact-free 3D sinogram, where $N\_{views}$ is the number of projection views, $U$ is the number of horizontal detector positions, and $V$ is the number of vertical detector rows. The motion-corrupted sinogram $S'$ is generated by applying a sequence of view-dependent geometric transformations to a copy of the original sinogram \cite{thies2025_cbctmotion}.

\paragraph{Motion Event Modeling.}
Patient motion is modeled as a small number of discrete inter-projection events, each representing a sudden positional change after which the object remains in its new configuration for all subsequent views \cite{thies2025_cbctmotion}. For each sinogram, the number of motion events  $N_e \in \{1, 2\}$ is randomly sampled. The onset position of each event is drawn uniformly from the range $[100,\, N\_{views}-100]$ to avoid boundary effects at the start and end of the acquisition. To simulate realistic motion trajectories, motion parameters are applied cumulatively across consecutive projections, resulting in temporally consistent motion patterns. 
\begin{itemize}
\item \textbf{In-plane rotation:} A rotation angle $\alpha_i \sim \mathcal{U}(-5^\circ, +5^\circ)$ is applied to all projections from the event onset onward \cite{thies2025_cbctmotion}. The rotation operates in the $(U, V)$ detector plane using cubic interpolation with constant zero-padding, preserving the original detector dimensions.

\item \textbf{Horizontal translation:} A shift $\Delta u_i \sim \mathcal{U}(-15, +15)$ pixels along the detector $U$-axis is applied starting 5 to 19 views after the rotation onset \cite{thies2025_cbctmotion}. This delay reflects the observation that rotational and translational components of real patient motion do not typically begin simultaneously. The shift is implemented via coordinate remapping with cubic interpolation and constant zero-padding \cite{Gonzalez2008}.
\end{itemize}
All rotations are applied first, followed by all translations. Since each transformation is applied from its event onset to all subsequent views, the effects accumulate across events, resulting in progressively increasing distortions. The use of constant zero-padding at detector boundaries introduces edge discontinuities in the transformed projections, which contribute to more pronounced and realistic artifacts in the reconstructed images \cite{wagner2003quantification}.
The ranges of translation and rotation are chosen to simulate realistic patient head motion observed during CT acquisition. 

The output sinogram is stored in 32-bit floating-point precision. The artifact generation supports incremental processing: sinograms for which an artifact file already exists are automatically skipped, and non-three-dimensional inputs are excluded.
The motion simulation parameters are summarized in Table \ref{tab:motion_parameters}. 

\begin{table}[H]
\centering
\caption{Motion simulation parameters for artifact generation}
\label{tab:motion_parameters}
\begin{tabular}{p{5cm} p{9cm}}
\hline
\textbf{Parameter} & \textbf{Value} \\
\hline
Motion type & Inter-projection rotation and translation \\
Number of motion events & 1--2 (randomly sampled) \\
Event onset range & $[100,\, N\_{views}-100]$\\
Rotation angle range &  $[-5^\circ, +5^\circ]$\\
Translation range & $\pm 15$ pixels (detector $U$ axis)\\
Translation delay & 5--19 views after rotation onset \\
Interpolation method & Cubic B-spline (order 3)\\
Boundary handling & Constant padding (value 0) \\
Application domain & 3D sinogram domain (per-view transformation) \\
Persistence & Each event affects all subsequent views \\
Output precision & 32-bit floating point \\
Motion severity & Moderate 1–2 motion events with bounded rotation and translation ranges\\
\hline
\end{tabular}
\end{table}

\paragraph{Discussion.}
Since motion artifacts are introduced in the projection-domain, the resulting inconsistencies propagate through the reconstruction process and manifest as structured artifacts such as streaking and blurring in the image-domain. This ensures that both pipelines operate on physically consistent artifact patterns. While more advanced models could incorporate non-rigid deformations or continuous trajectories, the current approach provides an effective and computationally efficient approximation suitable for generating training data. 


\section{Projection-Domain-Specific Data Processing }
\subsection{Conversion of 3D Sinograms to 2D Representation  }
Following motion artifact simulation, both clean and motion-corrupted sinograms are represented as 3D arrays of shape  $(N\_{views}, U, V)$, where $N\_{views}$  denotes the number of projection views, and $(U, V)$ represent the detector dimensions.

To enable efficient training, the 3D sinograms are decomposed into 2D slices along the detector row dimension.  Since there are 53 patients, a total of $53 \times  detector\_v$ = $53 \times 800$ = 42,400 2D slices are obtained for clean and motion-corrupted data. 

This slicing procedure is applied consistently to both clean and motion-corrupted sinograms, ensuring spatial alignment and forming paired training samples for supervised learning.

To further improve data quality, optional filtering is applied to exclude slices with negligible signal content. In particular, slices with mean intensity below a predefined threshold (\textbf{$\textbf{0.01}$}) are discarded. For the dataset considered, this results in approximately  16,902 slices after filtering.

\subsection{Sampling, Normalization, and Dataset Splitting}
\paragraph{Sampling}
To control dataset size and ensure computational efficiency, a maximum number of projection views per patient is specified as a configurable parameter.  When the number of available views exceeds this limit, a \textbf{random sampling strategy} is applied to select a subset of views. This introduces variability in the training data while maintaining a manageable number of samples per patient. In the experiments conducted in this work, the maximum number of views per patient is capped at \textbf{100} by command-line argument. 
\paragraph{Splitting}
Following dataset construction and per-patient view sampling, the data is split at the \textbf{patient level} into training (80\%), validation (10\%), and test (10\%) subsets. Patient identifiers are extracted from filenames, and patients are randomly shuffled before being assigned to each subset. All projection views corresponding to a given patient are exclusively assigned to a single subset, ensuring no data leakage across splits. In this work, the dataset is divided into \textbf{42 patients for training, 5 for validation, and 6 for testing},resulting in around 4100 training sinogram slices, depending on the number of available views per patient after sampling. 
\paragraph{Normalization}

During data loading, each sinogram is normalized independently using the intensity statistics of the artifact-corrupted input on a per-sample basis. Specifically, min–max normalization is applied using the minimum and maximum values of the corrupted sinogram, and the same scaling is applied to the corresponding clean target. This ensures a consistent intensity range across inputs and targets while preserving relative structural variations.

 

\section{Image-Domain-Specific Data Processing}
\subsection{Reconstruction to Image space}
Following motion artifact simulation which results in 3D sinogram pairs, reconstruction is performed using a filtered backprojection pipeline (FDK) \cite{feldkamp1984practical}, where frequency-domain filtering is applied prior to backprojection. \cite{diffct2025}. The reconstruction process has three main stages: geometric weighting, frequency-domain filtering, and backprojection. 

Prior to filtering, the projection data is weighted to account for cone-beam geometry. The weighting is computed based on detector coordinates and system geometry, where each detector element is scaled according to its distance from the X-ray source. Specifically, a distance-dependent weighting factor is applied to compensate for beam divergence, ensuring that rays farther from the central axis contribute appropriately during reconstruction. Next, a ramp filter combined with a Hann window is applied in the frequency domain to enhance high-frequency components while reducing noise. Finally, the filtered projections are backprojected across all acquisition angles to reconstruct the volumetric image. The same acquisition geometry used during forward projection is employed during reconstruction to ensure consistency between simulated projections and reconstructed volumes. 

Following reconstruction, the resulting CT volumes are normalized to ensure consistent intensity scaling across samples. Specifically, intensity values are clipped between the 1st and 99th percentiles to reduce the influence of extreme outliers, and subsequently rescaled to the range [0, 1]. 


\subsection{3D CT Volumes to 2D CT Slice Representation}

To enable efficient training, the reconstructed volumes are decomposed into axial 2D slices, where each slice is treated as an independent training sample. Two-dimensional slices are extracted along the axial direction of the reconstructed volume, corresponding to the depth axis of the 3D data. This slicing strategy aligns with standard clinical CT viewing planes and preserves consistent anatomical structures across slices, enabling the model to learn spatially meaningful artifact patterns. 

Additionally, slices with low structural information ($<0.01$) are excluded based on intensity variance, ensuring that the model focuses on informative anatomical regions. This gives a pairing of 10,000 clean-corrupted CT slices. 

\subsection{Slice Sampling, Normalization, and Dataset Splitting }
\paragraph{Sampling}
Following slice filtering, a controlled sampling strategy is applied to ensure balanced representation across patients. For each CT volume, a maximum of \textbf{140 slices per patient} is retained. When the number of valid slices exceeds this limit, a \textbf{uniform sampling strategy} is employed, selecting slices at regular intervals across the volume. The selected slices from all patients are then aggregated to form the complete dataset.
\paragraph{Splitting}
After sampling, all selected slices are aggregated and subsequently split at the \textbf{patient level} into training (80\%), validation (10\%), and test (10\%) subsets.  All slices from a given patient are assigned exclusively to one subset to prevent data leakage.  In this work, the dataset is divided into \textbf{42 patients for training, 5 for validation, and 6 for testing}, resulting in  approximately \textbf{5300 training slices}, with validation and test samples derived from their respective patient groups.
\paragraph{Normalization}
During data loading, each slice is normalized independently using the intensity statistics of the artifact-corrupted input. Min–max normalization is applied using the minimum and maximum values of the corrupted slice, and the same scaling is applied to the corresponding clean target. This normalization is performed on-the-fly during data retrieval, ensuring consistency across training, validation, and testing phases.

\begin{equation}
\begin{aligned}
S'_{\mathrm{norm}} &= \frac{S' - \min(S')}{\max(S') - \min(S') + \epsilon}, \\
S_{\mathrm{norm}}  &= \frac{S - \min(S')}{\max(S') - \min(S') + \epsilon}
\end{aligned}
\end{equation}

where $\epsilon$ is a small constant for numerical stability. This consistent normalization ensures that both input and target share the same intensity scale. Here $\epsilon$ is $1 \times 10^{-12}$. This normalization is applied independently of the initial preprocessing normalization and is specifically designed to stabilize reconstructed image intensities for learning. 

\section{Training Strategy}

Both projection-domain and image-domain pipelines are trained using supervised learning with paired data. In the projection-domain setting, models learn a mapping from motion-corrupted sinograms to clean sinograms, whereas in the image-domain setting, models learn a mapping from artifact-corrupted reconstructed image slices to clean reconstructed image slices.

 \subsection{Unified Training Framework  }
 A unified training framework is adopted to enable a controlled comparison between projection-domain and image-domain restoration strategies. All models, including U-Net, MR-LKV, SwinIR, Restormer, and RepLKNet, are trained independently in both domains using identical optimization settings and training procedures.

The key difference lies only in the data representation: projection-domain models operate on sinogram slices, while image-domain models operate on reconstructed CT slices. 


\subsection{Training Configuration}


All models are trained using the Adam optimizer with weight decay of $ 10^{-5}$. The learning rate is initialized based on the model type and dynamically adjusted using a \textbf{ReduceLROnPlateau } scheduler based on validation loss, with a reduction factor of 0.5 and a patience of 5 epochs.

The optimization objective is the L1 loss between the predicted and ground-truth outputs. In the projection-domain, this corresponds to sinogram slices, while in the image-domain, it corresponds to reconstructed CT image slices, as defined in Chapter \ref{chap:theoretical_background}. 




\subsection{Training Procedure and Stability Measures}
The training procedure remains consistent across both domains, with models trained using paired input-target samples 

To improve training stability, gradient clipping with a maximum norm in the range of $0.2 - 0.5$ is applied. Early stopping with a patience of 10 epochs is used to prevent overfitting. If required, model outputs are resized using bilinear interpolation to match ground truth dimensions before loss computation. 
Due to architectural constraints, input dimensions are adjusted for certain models. In particular, transformer-based models and large-kernel networks require input sizes divisible by specific factors, and inputs are cropped accordingly during training. Model performance is evaluated using Peak Signal-to-Noise Ratio (PSNR) and Structural Similarity Index (SSIM) on the validation and test sets. 

Despite the similarity in training procedure, the learning objectives differ fundamentally: projection-domain training focuses on correcting measurement inconsistencies, whereas image-domain training focuses on suppressing visually observable artifacts after reconstruction. 

\subsection{Validation and Model Selection}

After each epoch, the model is evaluated on the validation set. In addition to the loss, image quality metrics, including Peak Signal-to-Noise Ratio (PSNR) and Structural Similarity Index (SSIM) are computed to verify performance.

 The model achieving the lowest validation loss is saved as the best-performing model. Early stopping is employed to terminate training if no improvement in validation loss is observed for a fixed number of epochs, thereby preventing overfitting.

 \subsection{Testing and  Evaluation}

After training, the best model is evaluated on the unseen test set. Quantitative evaluation is performed using L1 loss, PSNR, and SSIM, providing a comprehensive assessment of both pixel-wise accuracy and structural fidelity. This evaluation is performed only after training to provide an unbiased assessment of model performance. 

 \subsection{Implementation Details}

All experiments are implemented using PyTorch and executed on a high-performance computing (HPC) cluster equipped with NVIDIA Tesla V100 GPUs with 32 GB memory. Data loading is performed using parallel workers to improve efficiency. Checkpoints are periodically saved to enable training resumption and reproducibility.  Model checkpoints, training logs, and validation curves are automatically saved during training to enable experiment tracking and result verification. 
To ensure reproducibility of the experiments, random seeds were fixed for all major libraries used in the training pipeline. Specifically, deterministic seeds were set for PyTorch, NumPy, and Python’s built-in random module. These settings control dataset shuffling, parameter initialization, and projection view sampling, ensuring reproducible results under identical configurations. 

Overall, the training strategy is designed to systematically evaluate the effectiveness of projection-domain and image-domain learning under controlled conditions. By training multiple architectures across both domains with consistent configurations, this work provides a comprehensive comparison of how domain choice influences motion artifact correction performance.

A summary of all training hyperparameters is provided in Table \ref{tab:training_params}. 

\begin{table}[H]
\centering
\caption{Training hyperparameters used for network optimization}
\label{tab:training_params}
\renewcommand{\arraystretch}{1.2}
\begin{tabular}{p{5cm} p{9cm}}
\toprule
\textbf{Parameter} & \textbf{Value} \\
\midrule
Optimizer & Adam \\
Learning rate & $2.5 \times 10^{-5}$ (MR-LKV, U-Net), $2 \times 10^{-5}$ (SwinIR, Restormer), $1 \times 10^{-5}$ (RepLKNet)\\
Batch size & 8 (MR-LKV, U-Net, RepLKNet), 2 (SwinIR), 1 (Restormer)\\
Epochs & 100 \\
Weight decay & $1 \times 10^{-5}$ \\
Scheduler & ReduceLROnPlateau\\
LR reduction factor & 0.5 \\
Scheduler patience & 5 epochs \\
Minimum learning rate & $1 \times 10^{-6}$ \\
Early stopping patience & 10 epochs \\
Max views per patient & 100 (per patient)\\
Patch size & Full resolution (MR-LKV, U-Net, RepLKNet); resized/cropped inputs (e.g., $128 \times 128$) for transformer-based models (SwinIR, Restormer)\\
Loss function & L1 loss \\
Gradient clipping & 0.2 - 0.5 (depending on domain)\\
Normalization & Min--max (artifact-based, per-sample)\\
\bottomrule
\end{tabular}
\end{table}



\section{Inference Pipeline}

\subsection{Projection-Domain Inference Pipeline}
\subsubsection{Dataset Preparation}
After training, the trained models are evaluated on unseen patient data. A separate test dataset consisting of 30 patients is processed using the same preprocessing pipeline as the training data, including DICOM-to-sinogram conversion and motion artifact simulation. 

The resulting 3D sinograms are converted into two-dimensional slices in the same manner as the training procedure. However, unlike the training pipeline, no slice filtering is applied at test time to preserve all projection information required for accurate volumetric reconstruction. As a result, a total of $30 \times 800$ = 24000 2D motion-corrupted sinogram slices are prepared.

\subsubsection{Generating Predicted 2D Sinogram}
Each 2D motion-corrupted sinogram of shape $(N\_{views}, U)$ is then passed through the trained network to obtain the corresponding restored sinogram. During inference, the input slices are normalized using the minimum and maximum values of the corrupted input slice. To satisfy architectural constraints, input slices are padded to dimensions divisible by model-specific factors (e.g., multiples of 8 or 16). After prediction, the output is clamped to the range $[0,1]$ and cropped back to the original spatial dimensions.

The predicted slices are then denormalized using the stored minimum and maximum values of the corresponding corrupted input. These normalization parameters are saved during inference to ensure consistent reconstruction.

Along with prediction, computational performance metrics, including model parameters, floating-point operations (FLOPs), and average inference time, are recorded for each model, providing information about efficiency and scalability.

\subsubsection{2D-to-3D Sinogram Aggregation}


Since models operate on individual 2D sinograms, a post-processing step is required to reconstruct full 3D volumes.

For each patient, the predicted 2D sinogram slices are aggregated by stacking them along the detector row dimension, restoring the original 3D sinogram structure. 

Metadata, such as normalization parameters, is saved during this process to ensure consistency in the reconstruction stage.

\subsubsection{CT Reconstruction from Predicted Sinograms  }



The reconstructed three-dimensional sinograms are subsequently converted into CT volumes using cone-beam geometry, implemented using the DiffCT framework \cite{diffct2025}.

Prior to backprojection, the sinograms undergo cosine weighting to account for cone-beam divergence, followed by filtering using a Hann-windowed ramp filter to suppress high-frequency noise and reduce ring artifacts. The Hann window is chosen for its favorable trade-off between spatial resolution and noise suppression.  The filtered projections are then backprojected over all acquisition angles to obtain volumetric reconstructions. 

The reconstruction process incorporates patient-specific metadata, including voxel spacing and volume dimensions, ensuring accurate spatial representation of anatomical structures. 
Following backprojection, the reconstructed volume is scaled and constrained to non-negative values to ensure physically meaningful attenuation coefficients.
The reconstructed CT volumes are further normalized and stored for subsequent evaluation and visualization.


This reconstruction step is essential, as it allows the performance of the proposed method to be assessed in the clinically relevant image-domain rather than solely in the projection-domain.

\subsection{Image-Domain Inference Pipeline}
\subsubsection{Dataset Preparation}
In the image-domain pipeline, the test dataset consists of reconstructed CT volumes obtained from motion-corrupted sinograms \cite{diffct2025}. For each patient in the test set, both artifact-corrupted and corresponding clean reconstructed volumes are available.
The reconstructed volumes are stored as three-dimensional arrays, where each volume consists of multiple axial slices. During inference, these volumes are loaded and processed sequentially. The input to the model is the artifact-corrupted reconstructed volume, while the clean volume is used as the ground truth for evaluation.
To ensure consistency with the training process, the same preprocessing steps, including normalization and data formatting, are applied during inference.

\subsubsection{Generation of Corrected 3D CT Volume  }
In the image-domain pipeline, inference is performed directly on reconstructed CT volumes to generate artifact-corrected outputs. Each motion-corrupted three-dimensional volume is processed in a slice-wise manner by decomposing it into two-dimensional axial slices.

Each slice is independently normalized on a per-sample basis using the intensity statistics of the artifact-corrupted slice to ensure compatibility with the training distribution. The normalized slice is then passed through the trained model to generate a corrected prediction. To accommodate architectural constraints of different models, input slices are padded where necessary, and the predicted outputs are cropped back to their original dimensions. The corresponding clean slice is normalized using the same scaling and clipped to the range [0, 1]. 

The predicted output is first rescaled back to the original intensity range using the inverse transformation. The corrected slices are then aggregated in their original spatial order to reconstruct the full three-dimensional CT volume. This process preserves the volumetric structure of the data while incorporating learned corrections at the slice level.

Unlike the projection-domain pipeline, no additional reconstruction step is required after inference, as the model operates directly on reconstructed image data. The final output is a corrected CT volume with reduced motion artifacts, which is used for quantitative evaluation and qualitative analysis.

During inference, the model efficiency and scalablity are also measured using computational performance metrics,like model parameters, floating-point operations (FLOPs), and average inference time. 








\section{Evaluation Strategy}

The performance of all models is evaluated in the reconstructed image-domain by comparing predicted volumes with corresponding ground truth volumes. This unified evaluation strategy is adopted to ensure a fair and clinically meaningful comparison between projection-domain and image-domain approaches.

In the projection-domain pipeline, the model produces corrected sinograms, which are subsequently reconstructed into CT volumes before evaluation. In contrast, in the image-domain pipeline, the model directly outputs corrected reconstructed volumes. Despite these differences, evaluation is consistently performed in the image-domain, as this reflects the final clinically relevant outcome.

During evaluation, both predicted and ground truth volumes are normalized using the intensity range of the clean reference volume to ensure consistent and fair comparison across models. The predicted volume is additionally clipped to the range [0, 1] to ensure numerical stability. Evaluation is performed on a per-slice basis and aggregated across the entire volume.

The following metrics are used:
\begin{itemize}

\item \textbf{PSNR (Peak Signal-to-Noise Ratio):} Measures pixel-wise reconstruction fidelity.

\item \textbf{SSIM (Structural Similarity Index):} Evaluates structural similarity between images.

\item \textbf{MAE (Mean Absolute Error):} Quantifies absolute intensity differences.

\item \textbf{RMSE (Root Mean Squared Error):} Measures overall reconstruction error.

\item \textbf{VIF (Visual Information Fidelity):} Assesses perceptual information preservation.

\item \textbf{LPIPS (Learned Perceptual Image Patch Similarity):} Evaluates perceptual similarity using deep features.

\end{itemize}

PSNR, SSIM, and VIF are computed on individual axial slices and averaged across all valid slices within each volume, while MAE and RMSE are computed over the entire normalized volume. LPIPS is computed on a subset of slices (every fifth slice) to reduce computational cost.  Slices with negligible signal content are excluded from evaluation to avoid bias in metric computation.

For each model, evaluation metrics are computed across all test patients and reported as mean and standard deviation. 








\section{Summary}
This chapter presents the complete methodology for motion artifact reduction in CT imaging, with a primary focus on evaluating the effectiveness of super-large-kernel convolution through the proposed MR-LKV model. Starting from DICOM data, sinograms are generated using cone-beam forward projection \cite{diffct2025}, followed by the simulation of motion artifacts directly in the projection-domain. To enable supervised learning, the generated 3D data is transformed into two-dimensional representations, allowing efficient training of deep learning models. The proposed MR-LKV architecture, designed with super-large receptive fields, is employed to capture the global spatial dependencies characteristic of motion artifacts, such as long-range streaking and structural distortions.

To comprehensively assess its effectiveness, the model is trained and evaluated in both projection-domain and image-domain settings. In the projection-domain pipeline, the model learns to correct motion-corrupted sinograms prior to reconstruction, while in the image-domain pipeline, it operates directly on reconstructed CT images to suppress visible artifacts.

During inference, predictions are aggregated to reconstruct three-dimensional volumes, which are evaluated in the reconstructed image-domain using quantitative and perceptual metrics. This unified evaluation framework ensures a fair and clinically relevant comparison of model performance.

It should be noted that the motion simulation process is synthetic and may not fully capture the complexity of real patient motion, which represents a limitation of the current study.

 
   % (\chapter{})
\cleardoublepage
\chapter{Results and Discussion}
\label{chap:results}

\section{Evaluation Overview}
This chapter presents a comprehensive evaluation of the research conducted in the context of this thesis by applying the described materials and methods of Chapter \ref{chap:methodology}. The primary objective is to assess the effectiveness of super-large-kernel convolution in modeling and suppressing motion-induced artifacts.


The proposed model is also compared with other CNN models, large-kernel models and transformer models. All models are trained and evaluated under identical experimental settings to ensure a fair comparison. The evaluation is conducted on a held-out test set consisting of previously unseen patient data, enabling assessment of both reconstruction accuracy and generalization capability. 
Although projection-domain and image-domain approaches are trained independently, they are intrinsically related through the CT reconstruction process. For clarity, results are presented first in the projection-domain, followed by the image-domain.





 

 
\section{Evaluation Metrics}


To comprehensively evaluate model performance, a combination of pixel-wise, structural, perceptual, and computational metrics is employed. Pixel-level accuracy is measured using PSNR, RMSE, and MAE, while structural similarity is assessed using SSIM. Perceptual quality is evaluated using LPIPS and VIF, capturing high-level visual consistency beyond pixel-wise agreement.

In addition, computational efficiency is analyzed using model parameters, FLOPs, and  average inference time, providing insights into practical deployment feasibility. 
These metrics are defined in detail in Chapter \ref{chap:theoretical_background}. 

 

\section{Preliminary Analysis of MR-LKV }
This section presents an initial analysis of the proposed MR-LKV model in isolation, focusing on its training behavior, qualitative reconstruction performance, and quantitative metrics in both projection- and image-domains. 

\subsection{Projection-Domain Results}
\subsubsection{Training Behaviour}

The training and validation loss curves demonstrate stable and consistent convergence throughout the training process. As shown in Fig. \ref{fig:training_loss_curve_projection}, the training loss decreases rapidly during the initial epochs from approximately \textbf{0.030} to below \textbf{0.010}, followed by a gradual reduction, reaching approximately \textbf{0.0055} towards the end of training. This behavior indicates effective optimization and progressive refinement of the model parameters. 

The observed performance improvements can be explained by the fact that motion artifacts arise from inconsistencies in the projection data. By operating directly in the sinograms, the model learns to correct these inconsistencies before they influence the reconstruction outcome.
\begin{figure}[htpb] 
    \centering
    \includegraphics[width=1.0\linewidth]{projection_loss_curve.pdf}
    \caption{Loss Curves and Validation Metrics in Projection-Domain Training}
    \label{fig:training_loss_curve_projection}
\end{figure}

The validation loss exhibits a similar downward trend, decreasing from around \textbf{0.027} and stabilizing near \textbf{0.012–0.013} after approximately 40–50 epochs. While minor fluctuations are observed in the early stages of training, these variations reduce over time, suggesting improved generalization. Although a gap between training and validation loss is present, it remains stable and does not increase, indicating that the model does not suffer from significant overfitting despite its relatively high capacity.

In addition, the validation PSNR shows a steady increase from approximately \textbf{27 dB} to around \textbf{33–34 dB}, reflecting continuous improvement in reconstruction quality. The SSIM values also exhibit a gradual upward trend, stabilizing close to \textbf{0.90}, further confirming the preservation of structural information in the reconstructed outputs.

 




\subsubsection{Quantitative Performance}

The performance of the MR-LKV model on projection-domain is evaluated using multiple image quality metrics, capturing both pixel-level accuracy and perceptual similarity. The obtained results are summarized as follows:

\begin{itemize}
    \item \textbf{PSNR}: 37.94 ± 10.20 dB
    \item \textbf{SSIM}: 0.9024 ± 0.0474
    \item \textbf{MAE}: 0.01234 ± 0.00781
    \item \textbf{RMSE}: 0.03097 ± 0.01193
    \item \textbf{VIF}: 0.3880 ± 0.1279
    \item \textbf{LPIPS}: 0.1027 ± 0.0515
\end{itemize}

The high PSNR and SSIM values indicate strong reconstruction fidelity and preservation of structural information. The relatively low MAE and RMSE further confirm accurate pixel-wise restoration. This performance can be attributed to the large receptive field of the MR-LKV architecture, which enables the model to capture long-range dependencies across the sinogram.  

The LPIPS score suggests that perceptual similarity between the reconstructed and ground truth images is well maintained, while the VIF score reflects moderate preservation of visual information content. These results indicate that the model not only minimizes numerical error but also preserves meaningful structural patterns, which is critical for downstream reconstruction quality.  

The relatively high standard deviation in PSNR indicates variability in performance across different slices, likely due to varying degrees of motion corruption.


\subsubsection{Computational Performance}
In addition to reconstruction quality, computational efficiency is an important consideration for practical deployment. 
The proposed model has approximately \textbf{8.06 million parameters} and requires \textbf{91.24 GFLOPs}, reflecting a moderate computational complexity relative to its performance.

The average inference time is measured to be \textbf{60.21 ms per sample}, indicating that the model is computationally efficient and suitable for near real-time applications, depending on deployment constraints. 

This balance between performance and complexity stems from the design of the MR-LKV architecture, which utilizes depthwise large-kernel convolutions to achieve a wide receptive field without significantly increasing the number of parameters. 

\subsubsection{Qualitative Analysis}
In the projection-domain, motion corruption introduces visible inconsistencies in the sinogram, appearing as discontinuities and irregular intensity variations across projection views, as shown in Fig. \ref{fig:sinogram_mrlkv}. These disruptions break the smooth sinusoidal patterns representing consistent object geometry, indicating a loss of global coherence in the acquired data.

\begin{figure}
    \centering
    \includegraphics[width=1.0\linewidth]{sinogram comparison.pdf}
    \caption{Qualitative comparison of clean, motion-corrupted, and proposed model 3D sinogram for projection-domain. }
    \label{fig:sinogram_mrlkv}
\end{figure}

The MR-LKV model effectively reduces these artifacts, producing visibly smoother and more coherent reconstructions. This improvement arises from the model's large receptive field, which enables it to capture globally distributed inconsistencies introduced by motion and promote spatial coherence throughout the image.  

The benefits of these corrections are clearly reflected in the reconstructed CT images as shown in Fig. \ref{fig:mrlkv_qualitative_projection}. 

\begin{figure}[htpb]
    \centering
    \includegraphics[width=1.0\linewidth]{CQ500CT19_mrlkv.png}
    \caption{MPR-based qualitative comparison of ground truth, motion-corrupted, and MR-LKV restored CT images with ROI zoom. }
    \label{fig:mrlkv_qualitative_projection}
\end{figure}
The MR‑LKV results show noticeable suppression of streaking artifacts, smoother intensity transitions, and enhanced structural continuity across all anatomical planes.

In axial slices, artifact suppression is most apparent along the cranial boundary, where edge definition in the highlighted regions appears sharper and more distinct. In coronal and sagittal views, where motion artifacts often manifest as inter‑slice inconsistencies, the reconstructions exhibit improved anatomical coherence and better preservation of overall structure. The ROI regions further emphasize reduced blurring and superior retention of fine details compared to the corrupted inputs.

Although variations in sinogram patterns may appear subtle, minor inconsistencies in projection space are magnified through the reconstruction process. As a result, even modest improvements in sinogram continuity lead to significant reductions in image‑domain artifacts.

However, some residual blurring and slight loss of sharpness persist near high‑frequency edges, particularly in heavily corrupted regions. This suggests that while MR‑LKV successfully captures global artifact structures, recovering extremely fine details remains challenging.


\subsection{Image-Domain Results}

\subsubsection{Training Behaviour}

The training and validation curves of the MR-LKV model in the image-domain show stable and effective learning behavior.   As shown in Fig. \ref{fig:training_loss_curve_image}, the training loss decreases steadily from roughly \textbf{0.045} to \textbf{0.022}, indicating efficient optimization.



\begin{figure}[htpb]
        \centering
        \includegraphics[width=1.0\linewidth]{mr_lkv_training_curves_images.png}
        \caption{Loss curves and Validation Metrics in Image-Domain Training}
        \label{fig:training_loss_curve_image}
    \end{figure}

The validation loss follows a gradual downward trend, stabilizing around 0.051 after approximately 30–40 epochs without any significant increase, indicating the absence of severe overfitting. The persistent gap between training and validation loss remains noticeable but stable, which is expected given the ill-posed nature of image-domain artifact correction, where motion artifacts are closely intertwined with anatomical features.

Compared to the projection-domain, learning in the image-domain poses a greater challenge since motion artifacts appear as complex spatial distortions and blurring rather than structured inconsistencies. Regardless, the validation PSNR increases from about $21.5\,\mathrm{dB}$ to $23.5\,\mathrm{dB}$, and SSIM reaches $0.86$--$0.87$, demonstrating effective reconstruction and good structural preservation. This performance highlights the advantage of MR-LKV's large receptive field, which allows the model to capture globally distributed artifact patterns while maintaining consistency across anatomical structures.



\subsubsection{Quantitative Performance}

The performance of the MR-LKV model on unseen test data is evaluated using multiple image quality metrics, capturing both pixel-level accuracy and perceptual similarity. The obtained results are summarized as follows:

\begin{itemize}
    \item \textbf{PSNR}: 29.48 ± 6.99 dB
    \item \textbf{SSIM}: 0.8904 ± 0.0556
    \item \textbf{MAE}: 0.03366 ± 0.01202
    \item \textbf{RMSE}: 0.05843 ± 0.01546
    \item \textbf{VIF}: 0.3540 ± 0.1250
    \item \textbf{LPIPS}: 0.1063 ± 0.0568
\end{itemize}

The PSNR and SSIM values indicate that the model achieves strong reconstruction fidelity  and preserves structural information to a decent extent. The relatively low MAE and RMSE confirm accurate pixel-wise restoration in image-domain. Together, these results highlight the model’s ability to strike a balance between artifact suppression and retention of anatomical structures. 

The LPIPS score reflects that perceptual similarity between the reconstructed and ground truth images is reasonably well maintained, while the VIF score indicates moderate preservation of visual information content.

However, compared to the projection-domain results, the PSNR and SSIM values in the image-domain are relatively lower. This is because image-domain correction operates on already reconstructed images, where motion artifacts have been propagated and spatially distributed during the reconstruction process. That makes it harder to remove the corruptions compared to the projection-domain, where the corruption is more directly represented in the sinogram space.

Despite this difficulty, MR-LKV continues to perform strongly owing to its capacity to capture long-range spatial dependencies. This allows the model to suppress large-scale artifact patterns effectively while preserving global anatomical consistency in the reconstructed images.

 

\subsubsection{Computational Performance}
The image-domain variant of the proposed model has approximately \textbf{7.69 million parameters} and requires \textbf{53.51 GFLOPs}, reflecting a computationally efficient design that still delivers a strong reconstruction performance. The average inference time is measured to be \textbf{28.66 ms per sample}, demonstrating improved computational efficiency and potential for near real-time deployment.  The convolutional architecture of MR-LKV facilitates efficient parallel computation, further reducing inference latency.


Compared to the projection-domain implementation, the image-domain model exhibits reduced computational overhead and faster inference, underscoring its scalability and practical advantage. This reduction results primarily from the tailored MR-LKV design for image-domain processing, where fewer channels and optimized kernel settings are sufficient to capture spatial dependencies in reconstructed images.

Overall, these results suggest that the image-domain MR-LKV model achieves a favorable trade-off between reconstruction quality and computational cost, making it well-suited for real-world and time-sensitive imaging scenarios.

\subsubsection{Qualitative Analysis}

Random slices (2D image) from all the three orientations- axial, sagittal, and coronal are chosen from a 3D input volume as showing in Fig.\ref{fig:mr_lkv_roi_image}
with corresponding zoomed-in regions of interest. 


\begin{figure}[htpb]
            \centering
            \includegraphics[width=1.0\linewidth]{CQ500CT169 CQ500CT169_mpr.png}
            \caption{MPR-based qualitative comparison of ground truth, motion-corrupted, and MR-LKV restored CT images with ROI zoom in image-domain}
            \label{fig:mr_lkv_roi_image}
        \end{figure}
The motion-corrupted slice images display noticeable motion-induced distortions, including blurring, streaking, and structural discontinuities, particularly visible along high-contrast boundaries such as the skull.
The MR-LKV reconstructed slice illustrates a clear indication of reduction in these artifacts across all anatomical planes. In the axial view, streaking patterns are effectively suppressed, transitions appear smoother, reflecting improved global consistency. The zoomed-in comparison along the ROIs further reveals improved edge definition and reduced blurring compared to the corrupted input.

In the coronal and sagittal views, where artifacts often appear as inter-slice inconsistencies, the motion-corrupted images reveal prominent distortions and discontinuities. The proposed method produces more coherent and anatomically consistent structures, indicating improved inter-slice continuity. This improvement results from the large receptive field of MR-LKV, enabling effective capture of long-range spatial dependencies and reinforcing structural consistency across slices.


Relative to the projection-domain displays, the image-domain reconstruction appears visually more refined, as evidenced by smoother textures and enhanced perceptual quality. This can be attributed to the model operating directly in the image space, where it learns to correct spatially distributed artifacts and optimize perceptual consistency. 

Minor residual blurring and slight sharpness loss are still noticeable along fine structural boundaries, particularly in high-frequency areas. This suggests that while MR-LKV effectively suppresses large-scale artifacts, recovering very fine details remains a challenge even in the image-domain.

Overall, the qualitative analysis demonstrates that the proposed super-large-kernel model balances artifact suppression and structural preservation, producing visually consistent and anatomically plausible reconstructions across multiple views. 




\section{Comparison of Models}
\subsection{Projection-Domain Comparison}
\subsubsection{Quantitative Results}
\begin{table}[H]
\centering
\caption{Quantitative comparison of reconstruction quality across different models in the projection-domain}
\label{tab:model_comparison_projection}
\renewcommand{\arraystretch}{1.2}

\resizebox{\textwidth}{!}{%
\begin{tabular}{lcccccc}
\toprule
\rowcolor{gray!20}
Model & PSNR $\uparrow$ & SSIM $\uparrow$ & VIF $\uparrow$ & RMSE $\downarrow$ & MAE $\downarrow$ & LPIPS $\downarrow$ \\
\midrule
U-Net & $34.89 \pm 10.44$ & $0.8664 \pm 0.0606$ & $0.3128 \pm 0.0954$ & $0.03840 \pm 0.01320$ & $0.01612 \pm 0.01031$ & $0.1400 \pm 0.0619$ \\
SwinIR &  34.40 ± 11.25&  0.8625 ± 0.0754&  0.3898 ± 0.1574&  0.04070 ± 0.01630&  0.01682 ± 0.01181&  0.1139 ± 0.0656\\
Restormer &  36.97 ± 10.01&  0.8991 ± 0.0456&  0.3725 ± 0.1107&  0.03315 ± 0.01114&  0.01322 ± 0.00806&  0.1101 ± 0.0563\\
RepLKNet & $35.33 \pm 10.50$ & $0.8786 \pm 0.0525$ & $0.3559 \pm 0.1002$ & $0.03829 \pm 0.01443$ & $0.01642 \pm 0.01111$ & $0.1405 \pm 0.0654$ \\
\rowcolor{green!10}
\textbf{MR-LKV (Proposed)} &  37.94 ± 10.20&  0.9024 ± 0.0474&  0.3880 ± 0.1279&   0.03097 ± 0.01193&  0.01234 ± 0.00781&  0.1027 ± 0.0515\\
\bottomrule
\end{tabular}
}
\end{table}


Table \ref{tab:model_comparison_projection} presents a comprehensive quantitative comparison of reconstruction quality across all evaluated models. The proposed MR-LKV model consistently achieves the best performance across most evaluation metrics.

In terms of PSNR, MR-LKV achieves the highest value (\textbf{37.94 dB}), indicating improved reconstruction fidelity compared to both conventional CNN-based models and transformer-based approaches. Similarly, MR-LKV attains the highest SSIM (\textbf{0.9024}), confirming improved structural preservation compared to other methods.

Error-based metrics further support these findings. MR-LKV achieves the lowest RMSE (\textbf{0.03097}) and MAE (\textbf{0.01234}), indicating reduced pixel-wise reconstruction error and improved numerical accuracy. These improvements suggest that the super-large-kernel design enables effective modeling of globally distributed inconsistencies introduced by motion artifacts.

In terms of perceptual quality, MR-LKV achieves the lowest LPIPS score (\textbf{0.1027}), reflecting improved perceptual similarity to ground truth images. While SwinIR achieves slightly higher VIF values (\textbf{0.3898}), the overall performance of MR-LKV remains competitive across both perceptual and structural metrics.

It is important to note that transformer-based models such as SwinIR and Restormer are inherently more expressive and may achieve higher performance under larger-scale training or higher computational budgets. However, under the constrained experimental setup considered in this work, MR-LKV demonstrates competitive performance while maintaining significantly lower computational requirements.



\subsubsection{Computational Efficiency Analysis }
\begin{table}[H]
\centering
\caption{Comparison of model efficiency and computational complexity in the projection-domain. FLOPs are computed as 2 × MACs using ptflops for all models.}
\label{tab:efficiency_comparison_-projection}
\begin{tabular}{lccc}
\toprule
\rowcolor{gray!20}
Model & Parameters (M) $\downarrow$ & FLOPs (G) $\downarrow$ & Inference Time (ms) $\downarrow$ \\
\midrule
U-Net &  7.76&  177.55&  \textbf{6.09}\\
SwinIR &  \textbf{1.04}&  860.52&  2514.62\\
Restormer &  10.66&   756.80&  227.13\\
RepLKNet &  2.73&  195.29&  193.40\\
\rowcolor{green!10} 
\textbf{MR-LKV (Proposed)} &  8.06&  \textbf{91.24}&  60.21\\
\bottomrule
\end{tabular}
\end{table}
The computational efficiency results summarized in the table \ref{tab:efficiency_comparison_-projection} reveal important trade-offs between model complexity, theoretical cost, and practical runtime performance across the evaluated architectures. 

The proposed MR-LKV achieves the lowest computational cost with \textbf{91.24 GFLOPs}, significantly reducing the number of operations compared to  conventional CNNs, RepLKNet and transformer-based models. This demonstrates that super-large-kernel depthwise convolution can efficiently capture global context without incurring excessive computational overhead.

In terms of runtime, MR-LKV achieves a balanced performance with an inference time of $\textbf{60.21 ms}$, markedly faster than transformer-based counterparts such as SwinIR and Restormer, while maintaining strong reconstruction quality. However, the simple architecture of U-Net makes its average Inference time ($\textbf{ 6.09 ms}$) much lesser than MR-LKV. 

Despite achieving a favorable balance between performance and efficiency, the inclusion of multiple architectural components introduces a higher parameter count ($\textbf{8.06M}$). As a result, the model may not be as lightweight as simpler CNN-based approaches or RepLKNet for highly resource-constrained deployments. 

Overall, these findings demonstrate that MR-LKV strikes a well‑balanced trade-off between accuracy and efficiency, making it a practical and scalable solution for projection-domain motion artifact correction.


\subsubsection{Trade off Between Performance and Efficiency}


To further analyze the relationship between reconstruction quality and computational cost, additional comparisons are presented through trade-off plots, such as FLOPs versus PSNR and FLOPs versus inference time.
\begin{figure}[htpb]
    \centering
    \includegraphics[width=1.0\linewidth]{projectiontradeoff.pdf}
    \caption{Trade-off between reconstruction performance and computational efficiency across different models, illustrated using FLOPs versus PSNR and inference time comparisons. }
    \label{fig:tradeoff_projection}
\end{figure}


As shown in Fig. \ref{fig:tradeoff_projection}, different models reveal distinct trade-offs between performance and efficiency. Transformer-based models such as SwinIR and Restormer achieve high reconstruction quality but at the cost of significantly higher computational complexity and inference time.

 In contrast, conventional CNNs such as U-Net and large-kernel RepLKNet offer lower computational demands but achieve limited performance 

The proposed MR-LKV model achieves an optimal balance in this trade-off space, combining high PSNR with notably lower FLOPs and faster inference compared to transformer-based methods. This indicates that MR-LKV effectively balances reconstruction quality and computational efficiency.

 In summary, these results demonstrate that super-large-kernel convolution enables efficient modeling of global dependencies without incurring the high computational overhead associated with self-attention mechanisms.



 

\subsubsection{Qualitative Results}
Qualitative comparisons in the projection-domain are illustrated in Fig. \ref{fig:sinogram_visual_comparison} and Fig. \ref{fig: cq500ct225}  which display both sinogram reconstructions and their corresponding CT images for all evaluated models.
\begin{figure}[htpb]
    \centering
    \includegraphics[width=1.0\linewidth]{sinogram comparison all.pdf}
    \caption{Qualitative comparison of clean, motion-corrupted, and model-predicted sinograms across different architectures. }
    \label{fig:sinogram_visual_comparison}
\end{figure}
\begin{figure}[htpb]
    \centering
    \includegraphics[width=1.0\linewidth]{CQ500CT225_roi.png}
    \caption{MPR visualization comparing motion artifact reduction across multiple models, including U-Net, RepLKNet, SwinIR, Restormer, and MR-LKV in projection-domain}
    \label{fig: cq500ct225}
\end{figure}

In the projection-domain, motion corruption introduces discontinuities and irregular intensity variations across projection views, disrupting the sinusoidal patterns and overall data consistency. 

 Conventional CNN models such as U-Net reduce local noise but fail to fully restore consistency across projections, leaving noticeable residual discontinuities. RepLKNet improves smoothness yet tends to over‑smooth projection patterns, leading to some loss of structural detail. Transformer-based models such as SwinIR and Restormer capture broader dependencies and produce more coherent sinograms, though they still show incomplete correction in regions with severe corruption.

 In contrast, the proposed MR-LKV model generates properly balanced, structurally consistent sinograms, effectively reducing discontinuities and improving continuity across projection angles. This performance can be credited to its use of super‑large-kernel convolutions, which allow direct modeling of long‑range relationships and enable the network to correct motion‑induced inconsistencies at their origin.

 While visual differences in sinograms may appear subtle, small inconsistencies in projection space can amplify through the reconstruction process. Consequently, even minor improvements in sinogram continuity lead to significant reductions in artifacts within the final CT images.

 This effect is evident in the corresponding CT reconstructions, where MR-LKV demonstrates superior suppression of streaking artifacts, enhanced structural continuity, and more precise preservation of anatomical boundaries compared with all baseline methods. ROI regions show clearer edge definition and reduced blurring, reflecting better recovery of fine details.

However, in severely corrupted regions, faint residual inconsistencies in the sinogram may still propagate into minor artifacts after reconstruction. Despite this, MR-LKV consistently achieves the best balance between artifact suppression and structural preservation among all evaluated methods.


 

 \subsubsection{Limitation}
Despite its strong performance, the proposed approach has a notable limitation.
It operates on 2D sinogram representations and does not explicitly exploit inter-slice consistency. Extending the model to 3D representations could further enhance performance. 

\subsection{Image-Domain Comparison}


\subsubsection{Quantitative Results}

The quantitative performance of different models in the image-domain setting is summarized in Table \ref{tab:model_comparison_image}. The evaluation is conducted on reconstructed CT volumes using multiple metrics, including PSNR, SSIM, RMSE, MAE, VIF, and LPIPS.
\begin{table}[H]
\centering
\caption{Quantitative comparison of reconstruction quality across different models in the image-domain}
\label{tab:model_comparison_image}
\renewcommand{\arraystretch}{1.2}

\resizebox{\textwidth}{!}{%
\begin{tabular}{lcccccc}
\toprule
\rowcolor{gray!20}
Model & PSNR (dB) $\uparrow$& SSIM $\uparrow$ & VIF $\uparrow$ & RMSE $\downarrow$ & MAE $\downarrow$ & LPIPS $\downarrow$ \\
\midrule
U-Net & 29.14 ± 7.06& 0.9002 ± 0.0510& 0.3668 ± 0.1528& 0.06183 ± 0.01842& 0.03315 ± 0.01303& 0.1059 ± 0.0544\\
SwinIR &  29.07 ± 7.22&  0.8907 ± 0.0604&  0.3630 ± 0.1250&  0.06268 ± 0.01789&  0.03613 ± 0.01308&  0.1008 ± 0.0576\\
Restormer &  29.80 ± 6.93&  0.9148 ± 0.0427&  0.3574 ± 0.1458&  0.05828 ± 0.01940&  0.02887 ± 0.01153&  0.0878 ± 0.0447\\
RepLKNet & 29.11 ± 7.10& 0.9020 ± 0.0440& 0.3222 ± 0.1136& 0.06156 ± 0.01601& 0.03133 ± 0.01038& 0.1090 ± 0.0545\\
\rowcolor{green!10}
\textbf{MR-LKV (Proposed)} &  29.48 ± 6.99 &  0.8904 ± 0.0556&  0.3540 ± 0.1250&   0.05843 ± 0.01546&  0.03366 ± 0.01202&  0.1063 ± 0.0568\\
\bottomrule
\end{tabular}
}
\end{table}
Among all models, Restormer achieves the highest PSNR (\textbf{29.80 dB}) and SSIM(\textbf{0.9148}), along with the lowest MAE (\textbf{0.02887}) and LPIPS (\textbf{0.0878}), indicating strong reconstruction fidelity and perceptual quality. This can be attributed to its transformer-based architecture, which effectively models global context through self-attention mechanisms.

The proposed MR-LKV model, while not achieving the highest SSIM or LPIPS values, demonstrates competitive performance across multiple metrics, with a PSNR of \textbf{29.48 dB} and RMSE of \textbf{0.05843}, closely approaching the performance of Restormer. This indicates that MR-LKV is able to achieve strong pixel-level reconstruction accuracy despite the increased complexity of image-domain artifact correction. However, its SSIM and LPIPS values are slightly lower compared to the best-performing model, indicating that while overall intensity reconstruction is accurate, the preservation of fine structural details and perceptual similarity remains comparatively limited.


The comparatively lower SSIM and perceptual metrics of all the models in the image-domain can be attributed to the inherent differences between projection and image spaces.

Despite these differences, MR-LKV maintains a strong balance between reconstruction accuracy and robustness, achieving competitive performance without the high computational overhead associated with transformer-based approaches.

 

\subsubsection{Computational Efficiency Analysis}

The computational efficiency of the models is evaluated in terms of parameter count, floating point operations (FLOPs), and average inference time per slice, as shown in Table \ref{tab:efficiency_comparison_image}

\begin{table}[H]
\centering
\caption{Comparison of model efficiency and computational complexity in the image-domain. FLOPs are computed as 2 × MACs using ptflops for all models.}
\label{tab:efficiency_comparison_image}
\begin{tabular}{lccc}
\toprule
\rowcolor{gray!20}
Model & Parameters (M) $\downarrow$ & FLOPs (G) $\downarrow$ & Inference Time (ms) $\downarrow$ \\
\midrule
U-Net &  7.76&  109.63&  \textbf{7.52}\\
SwinIR &  \textbf{1.04}&  175.34&  236.41\\
Restormer &  10.66&   456.94&  117.82\\
RepLKNet &  2.73&  118.69&  105.11\\
\rowcolor{green!10} 
\textbf{MR-LKV (Proposed)} &  7.69&  \textbf{53.51}&  28.66\\
\bottomrule
\end{tabular}
\end{table}

MR-LKV demonstrates a favorable efficiency profile, achieving the lowest computational complexity with only \textbf{53.51 GFLOPs}, significantly lower than all other models. In comparison, Restormer requires 456.94 GFLOPs, making it substantially more computationally intensive.

In terms of inference time, U-Net is the fastest (7.52 ms), followed by MR-LKV (\textbf{28.66 ms}), which remains significantly faster than transformer-based architectures such as SwinIR (236.41 ms) and Restormer (117.82 ms) or even another large-kernel CNN model, RepLKNet (105.11 ms).

Despite having a comparable number of parameters to U-Net (\textbf{7.7M}) and higher parameter count than RepLKNet (\textbf{2.73 M}), MR-LKV achieves improved computational efficiency. This can be attributed to its architectural design, where large receptive fields are achieved using depthwise separable convolutions and selectively applied large kernels, rather than repeatedly stacking computationally expensive operations across multiple layers. As a result, MR-LKV is able to capture long-range spatial dependencies efficiently without incurring significant computational overhead.

Compared to the projection-domain configuration, the image-domain model demonstrates lower computational cost and faster inference. This improvement arises from the more spatially structured nature of reconstructed images, enabling efficient artifact correction through fewer operations and a streamlined model design.


\subsubsection{Trade-off Between Performance and Efficiency}

As shown in Fig. \ref{fig:tradeoff_image}, the models exhibit clear trade-offs between reconstruction performance and computational efficiency in the image-domain.  
\begin{figure}[htpb]
    \centering
    \includegraphics[width=1.0\linewidth]{psnrtradeoffimage.pdf}
    \caption{Trade-off between reconstruction performance and computational efficiency across different models, illustrated using FLOPs versus PSNR and inference time comparisons for image-domain}
    \label{fig:tradeoff_image}
\end{figure}

The proposed MR-LKV model  has a favorable intermediate position, achieving competitive PSNR while maintaining relatively low computational complexity.  Despite having much fewer FLOPs than transformer-based models, MR-LKV delivers improved reconstruction quality compared to other lightweight baselines such as U-Net or RepLKNet (large-kernel), demonstrating the effectiveness of its super-large-kernel design in capturing global contextual information.

Even in terms of efficiency, MR-LKV maintains a favorable position. While models like SwinIR and Restormer incur higher inference time due to their architectural complexity, MR-LKV achieves significantly faster inference while also giving strong performance. In this regards, U-Net demonstrates the lowest inference time but at the cost of reduced reconstruction accuracy.

Overall, the MR-LKV model provides an effective balance between performance and efficiency, bridging the gap between lightweight CNN-based models and computationally intensive transformer-based approaches in the image-domain.

 

\subsubsection{Qualitative Results}

Qualitative comparisons of reconstructed CT images by using image-domain trained models are presented in Fig. \ref{fig:qual_comparison_image}.
The motion-corrupted inputs present pronounced motion-induced distortions, including blurring, streaking, and structural discontinuities, particularly along high-contrast boundaries.

\begin{figure}[htpb]
    \centering
    \includegraphics[width=1.0\linewidth]{CQ500CT243 CQ500CT243_mpr.png}
    \caption{MPR visualization comparing motion artifact reduction across multiple models, including U-Net, RepLKNet, SwinIR, Restormer, and MR-LKV in image-domain}
    \label{fig:qual_comparison_image}
\end{figure}
In the axial view, MR-LKV effectively suppresses motion-induced streaking while maintaining clearer boundaries compared to U-Net and RepLKNet. The zoomed regions show improved edge definition and reduced blurring, indicating better local feature recovery. In contrast, U-Net exhibits residual streaking and blurring, while RepLKNet tends to oversmooth fine details, leading to loss of edge sharpness.

In the coronal and sagittal views, MR-LKV shows a notable improvement in inter-slice consistency, producing more coherent anatomical structures compared to baseline CNN models. This improvement can be attributed to its large receptive field, which enables the model to capture long-range spatial dependencies and enforce structural consistency across slices more effectively than standard convolutional models.


However, in comparison to transformer-based models, MR-LKV exhibits slightly less aggressive smoothing of residual artifacts, with faint streaking patterns still visible in some regions. Additionally, in highly corrupted areas, the model may not fully recover fine-grained structures to the same extent as Restormer, which benefits from its higher capacity and global attention mechanisms.

It is also observed that reconstructions in the image-domain appear visually more refined than those obtained in the projection-domain. This can be attributed to the fact that, during reconstruction, motion artifacts become spatially distributed, enabling the model to directly learn and correct these structured distortions within the image space. Consequently, the results exhibit smoother textures and better perceptual quality, even though the underlying correction task remains inherently more complex.




\section{Discussion}

The results across both projection-domain and image-domain settings demonstrate that super-large-kernel convolution provides an effective alternative to both conventional CNNs and transformer-based architectures for CT motion artifact reduction. 

The effectiveness of MR-LKV can be attributed to its super-large-kernel design, which allows direct modeling of long-range dependencies. Unlike conventional CNNs with limited receptive fields, MR-LKV captures global context more effectively, while remaining computationally efficient due to its convolutional structure. Compared to transformer-based models, which rely on self-attention for global modeling, MR-LKV provides a more efficient alternative with significantly lower computational cost.

In the projection-domain, MR-LKV shows strong  quantitative performance by directly addressing inconsistencies at their source in the sinogram space. The inclusion of frequency-enhanced processing and view-axis attention further improves the model’s ability to capture structured distortions across projection views. 
In the image-domain, although the model operates on already reconstructed data where artifacts have propagated, MR-LKV continues to provide effective artifact suppression through super-large-kernel spatial modeling. While the quantitative performance is slightly lower compared to the projection-domain approach, the image-domain model offers a more computationally efficient, visually corrected and simplified pipeline, making it suitable for scenarios where direct access to projection data is not available.

The comparison with transformer-based models highlights an important trade-off. While transformer architectures achieve strong global modeling through self-attention, they incur significantly higher computational cost. In contrast, MR-LKV achieves competitive performance using convolutional operations with large receptive fields, providing a more efficient alternative for modeling global dependencies in high-resolution medical imaging tasks.


Overall, the proposed approach successfully bridges the gap between lightweight CNNs and computationally expensive transformers, providing a balanced trade-off between reconstruction quality and efficiency.
This validates the central hypothesis of the thesis and demonstrates the practical applicability of super-large-kernel CNNs for motion artifact reduction.







 





  % (\chapter{})
%% ... more chapters ....
\chapter{Conclusion}
\label{chap:conclusion}

\section{Summary of Work}
This thesis investigated the effectiveness of super-large-kernel CNN in motion artifact reduction in computed tomography (CT) imaging. Motion artifacts, arising from inconsistencies across projection views during data acquisition, are a major challenge due to their globally distributed nature and their impact on diagnostic image quality. Traditional convolutional neural networks (CNNs), which rely on small kernel sizes, often struggle to model such long-range dependencies efficiently, while transformer-based approaches, although effective, have high computational cost.

To address this gap, a novel architecture, MR-LKV (Motion Restoration using Large Kernel Convolution), is proposed. The model leverages super-large depthwise convolutional kernels to capture global spatial dependencies, enabling effective modeling of globally distributed motion artifacts. The architecture is designed to maintain computational efficiency through depthwise separable convolutions and a lightweight encoder–decoder framework.
A unified framework is developed, including DICOM preprocessing, sinogram generation via forward projection, motion artifact simulation, and training in both projection and image domains. In the projection domain, MR-LKV operates directly on sinogram data to correct motion-induced inconsistencies prior to reconstruction, while in the image domain, a simplified variant focuses on spatial refinement of reconstructed CT slices.


The proposed approach is compared against multiple baseline architectures, including U-Net, RepLKNet, SwinIR, and Restormer, using quantitative, perceptual, and computational efficiency metrics, along with qualitative analysis. It is demonstrated by experimental results that MR-LKV achieves competitive reconstruction quality while maintaining significantly lower computational complexity compared to transformer-based architectures.

 Overall, this work demonstrates that super-large-kernel convolutional architectures provide an effective and efficient solution for modeling global dependencies in CT artifact reduction, offering a practical alternative to both conventional CNNs and computationally intensive transformer-based models.  
 


\section{Main Findings}

The results of this work provide several important insights into the role of super-large kernel convolution for motion artifact reduction in CT imaging.

 First, a clear distinction is observed between projection-domain and image-domain correction. In the projection domain, MR-LKV achieves the best quantitative performance, effectively restoring consistency in the sinogram and reducing reconstruction errors. However, the visual improvements in this domain may appear subtle, as small inconsistencies in projection space are amplified during reconstruction. In contrast, the image domain presents a more challenging setting where artifacts are already spatially distributed, resulting in slightly lower quantitative metrics. Despite this, image-domain results appear visually more refined, as the model directly optimizes perceptual quality, leading to smoother textures and improved anatomical coherence.

 Second, the findings demonstrate that the architectural design of MR-LKV enables effective modeling of globally distributed artifacts. Compared to conventional CNN-based models such as U-Net and RepLKNet, MR-LKV consistently produces improved structural consistency and reduced artifact intensity. While transformer-based models such as Restormer achieve better perceptual metrics and finer detail recovery, they do so at the cost of significantly higher computational complexity 

From a computational standpoint, MR-LKV achieves the lowest FLOPs among all tested models while maintaining moderate parameter count and considerably faster inference compared to transformer-based architectures. This balance between performance and efficiency highlights its potential for practical deployment in clinical settings. 

Qualitative analysis further supports these findings, showing that MR-LKV effectively removes distributed motion artifacts globally while preserving fine anatomical details. In particular, ROI-based comparisons highlight improved edge sharpness and reduced streak artifacts compared to baseline models.

Overall, the results validate super-large kernel convolution as an effective and efficient approach for motion artifact reduction in CT imaging.

 

\section{Limitations}

Despite its competitive performance, MR-LKV does not consistently outperform transformer-based models such as Restormer in all evaluation metrics, particularly in terms of structural similarity (SSIM) and perceptual similarity (LPIPS). This may be attributed to the fact that transformer-based architectures are inherently well-suited for modeling long-range dependencies through self-attention mechanisms. Additionally, while the model achieves a favorable trade-off between performance and efficiency, the use of large kernels still introduces moderate computational overhead compared to lightweight CNN architectures.


Secondly, although the projection-domain approach achieves higher reconstruction accuracy, it requires raw sinogram data and has a very complex processing pipeline, which limits its usability in scenarios where only reconstructed images are available. Conversely, the image-domain approach, while more efficient and flexible, functions on already corrupted data, making complete artifact removal inherently more difficult. 


Thirdly, the evaluation is conducted on synthetically generated motion artifacts, which may not fully represent the complexity of real-world patient motion. This may limit the generalizability of the observed results. Additionally, the study focuses primarily on single-modality CT data and does not consider multi-modal or multi-task learning scenarios, which could further enhance robustness and performance. 

Finally, although MR-LKV achieves a favorable balance between performance and efficiency, there is still a trade-off between computational cost and reconstruction accuracy. Further optimization may be required for deployment in highly resource-constrained environments. 

 
 


\section{Future Work}

Future work should focus on evaluating the proposed approach on real clinical datasets containing naturally occurring motion artifacts to assess practical applicability.


Extending the framework to three-dimensional processing could further improve volumetric consistency by capturing inter-slice dependencies. Future work could explore hybrid architectures that combine super-large kernel convolution with more expressive attention mechanisms. Integrating adaptive self-attention with large receptive fields may further improve the modeling of complex and spatially varying motion artifacts, while maintaining computational efficiency.

Further optimization of baseline architectures under comparable computational settings would enable more comprehensive benchmarking. Finally, exploring model compression, pruning, and hardware-specific optimization could facilitate real-time deployment in clinical environments.

In addition, integrating projection-domain and image-domain correction within a unified end-to-end framework could leverage the strengths of both approaches, enabling more robust artifact reduction across the entire reconstruction pipeline.


Overall, this work demonstrates that super-large kernel convolution provides a robust and computationally efficient framework for modeling global motion artifacts, offering a promising direction for future research in CT image restoration.


 

 

 


     % (\chapter{})
%\cleardoublepage

\appendix
\cleardoublepage
\include{appendix}   % appendix A
\cleardoublepage
%\include{thesis10}   % appendix B
%\cleardoublepage
%\include{thesis11}   % appendix C
%\cleardoublepage

%% Do not change, auto-generated lists of figures, tables and literature %%
% Glossar
\printunsrtglossary[type=abbreviations]
\cleardoublepage

% List of figures
\addcontentsline{toc}{chapter}{\listfigurename}
\listoffigures
\cleardoublepage

% List of tables
\addcontentsline{toc}{chapter}{\listtablename}
\listoftables
\cleardoublepage


\chapter*{List of Mathematical Notations}
\addcontentsline{toc}{chapter}{List of Mathematical Notations}

\renewcommand{\arraystretch}{1.3}

\begin{longtable}{ll}

\textbf{Symbol} & \textbf{Description} \\
\hline
\endfirsthead

\textbf{Symbol} & \textbf{Description} \\
\hline
\endhead

\multicolumn{2}{l}{\textbf{CT Imaging and Physics}} \\
\hline
$I$ & Measured X-ray intensity \\
$I_0$ & Initial X-ray intensity \\
$\mu(x)$ & Linear attenuation coefficient at position $x$ \\
$f(x,y)$ & Attenuation distribution in spatial domain \\
$p(t,\theta)$ & Projection (sinogram) at detector position $t$ and angle $\theta$ \\
$\theta$ & Projection angle \\
$t$ & Detector position \\

\multicolumn{2}{l}{\textbf{Reconstruction and Data Representation}} \\
\hline
$S$ & Ground truth image or sinogram \\
$\hat{S}$ & Predicted (reconstructed) image or sinogram \\
$N$ & Number of pixels or elements \\
$x, y$ & Spatial coordinates \\
$A$ & Number of projection angles \\
$H, W$ & Height and width of image or sinogram \\

\multicolumn{2}{l}{\textbf{Deep Learning and CNN Operations}} \\
\hline
$x$ & Input image or feature map \\
$y$ & Output feature map \\
$w$ & Convolution kernel \\
$X$ & Input feature tensor \\
$Y$ & Output feature tensor \\
$C$ & Number of channels \\
$k_l$ & Kernel size at layer $l$ \\
$s_l$ & Stride at layer $l$ \\
$R_l$ & Receptive field at layer $l$ \\

\multicolumn{2}{l}{\textbf{MR-LKV Model Components}} \\
\hline
$X'$ & Feature map after pointwise ($1 \times 1$) convolution \\
$U$ & Intermediate feature map after depthwise convolution \\
$X_{ca}$ & Channel-attended feature representation \\
$X_{vaa}$ & View-axis attention feature representation \\
$w$ & Channel attention weights \\
$a$ & View-axis attention weights \\
$f_\theta$ & Neural network mapping function with parameters $\theta$ \\
$\theta$ & Learnable model parameters \\

\multicolumn{2}{l}{\textbf{Frequency Domain Representation}} \\
\hline
$\hat{X}$ & Frequency-domain representation of feature map \\
$\mathcal{F}$ & Fourier transform (FFT) \\
$\mathcal{F}^{-1}$ & Inverse Fourier transform (IFFT) \\


\end{longtable}

% Literature list
% %CONFIG:
%\selectlanguage{german}{\addcontentsline{toc}{chapter}{\bibname}}
\selectlanguage{english}{\addcontentsline{toc}{chapter}{\bibname}}

\printbibliography
\end{document}


